% concatenated from Cluster_ev-structure_of_earth_oscillation.tex

\documentclass[a4paper,12pt]{article}

\usepackage{amsmath,amsthm,amssymb}
\usepackage{enumerate}
\usepackage{graphicx}
\usepackage{titlesec}
\usepackage{empheq}
\usepackage{color}
\usepackage{xfrac,bigints}
\usepackage{comment}
\usepackage{subfigure}
\usepackage{caption}
\usepackage[normalem]{ulem}

% \usepackage[backend=biber,style=numeric]{biblatex}
% \addbibresource{references.bib}

\topmargin=-1.5cm
\textheight=23cm
\oddsidemargin=0cm
\textwidth=16cm

\newtheorem{Th}{Theorem}[section]
\newtheorem{lemma}[Th]{Lemma}
\newtheorem{remark}[Th]{Remark}
\newtheorem{Prop}[Th]{Proposition}
\newtheorem{Cor}[Th]{Corollary}
\newtheorem{Def}[Th]{Definition}
\newtheorem{Prob}[Th]{Problem}
\newtheorem{Conj}[Th]{Conjecture}
%\newtheorem{Rem}[Th]{Remark}

\newtheorem{a priori condition}[Th]{A priori condition}

\newtheorem{assumption}[Th]{Assumption}

\newtheorem{properties}[Th]{Properties of relaxation tensor}

\titleformat*{\section}{\normalsize\bfseries}
\titleformat*{\subsection}{\normalsize\bfseries}
\renewcommand{\thesection}{{{\normalsize\arabic{section}}}}
\renewcommand{\thesubsection}{{{\thesection.\normalsize\arabic{subsection}}}}
\renewcommand{\theequation}{{{\thesection.\arabic{equation}}}}
\renewcommand{\contentsname}{{\large Contents}}
\renewcommand{\refname}{{\large References}}

\newcommand{\argmax}{\mathop{\rm arg~max}\limits}
\newcommand{\argmin}{\mathop{\rm arg~min}\limits}

\newcommand{\eproof}{\mbox{\ }~\hfill
\mbox{\large $\Box$}}

\newcommand{\pf}{\noindent{\bf Proof}}

\DeclareMathOperator\supp{supp}

\DeclareMathOperator*{\esssup}{ess\,sup}


\def\Z{{\mathbb Z}}
\def\N{{\mathbb N}}
\def\R{{\mathbb R}}
\def\C{{\mathbb C}}
\def\DS{\displaystyle}
\def\vep{\varepsilon}
\def\sgn{\mbox{\rm sgn}}
\def\dvg{\mbox{\rm div}}

\def\red{\color{red}}
\def\blue{\color{blue}}

\definecolor{purple}{rgb}{.75, 0, .25}
\def\purple{\color{purple}}

\makeatletter
\long\def\@makefntext#1{\parindent 1em\noindent
\@hangfrom{\hbox to 1.8em{\hss$^{\@thefnmark}$}}#1}
\makeatother

\begin{document}

\title{Explicit form of relaxation tensor for isotropic extended Burgers model and its spectral inversion}

\author{
Youjun Deng\thanks{School of Mathematics and Statistics, Central South University,Changsha 410083, P.R. China (youjundeng@csu.edu.cn).}
\and
Ching-Lung Lin\thanks{Department of Mathematics, National Cheng-Kung University, Tainan 701, Taiwan (Email:
cllin2@mail.ncku.edu.tw).}
\and Gen Nakamura\thanks{Department of
Mathematics, Hokkaido University, Sapporo 060-0808, Japan and Research Center of Mathematics for Social Creativity, Research Institute for Electronic Science, Hokkaido University, Sapporo 060-0812, Japan (Email: gnaka@math.sci.hokudai.ac.jp).}}

\date{}
\maketitle

\begin{abstract}
Concerning the anelastic nature of Earth, the quasi-static extended Burgers model (abbreviated by q-EBM), an integro-differential system, is used to study the free oscillation of Earth (abbreviated by FOE). In this paper, we first provide a general method to obtain an explicit form of the relaxation tensor for inhomogeneous isotropic q-EBM. Then, we apply it to compute the eigenvalues of the free oscillation of Earth, assuming that Earth is a unit ball modeled as a homogeneous and isotropic q-EBM. So far, an analytical and systematic way to compute the eigenvalues of the FOE has been missing when modeling Earth as a q-EBM. In particular, we compute some clusters of eigenvalues (abbreviated by C-ev's). To be more precise, integrating by parts with respect to time of the q-EBM under the assumption that the initial strain is zero, the q-EBM becomes the sum of two terms. 
\begin{comment}of divergence form with negative definite second-order elliptic operators acting on the displacement vector. 
\end{comment}
The first term, called the instantaneous term, doesn't have any integration with respect to time, but the second term, called the memory term, has such an integration.  Then, consider the eigenvalues of the instantaneous part of the q-EBM. The eigenfunctions of C-ev's share the same eigenfunctions of the instantaneous part. However, the C-ev's may be shifted from the eigenvalues of the instantaneous part. Further, we analyze the structure of C-ev's and provide an inversion formula identifying the q-EBM from the C-ev's. 
%In addition, we provide an inversion method that recovers the relaxation tensor of the EBM by knowing the respective two sets of cluster v-ev and that of d-ev.


\bigskip

\noindent
{\bf Keywords:} anelastic, viscoelasticity, relaxation tensor, extended Burgers model, quasi-static Boltzmann-type viscoelastic system.

\noindent
{\bf MSC(2010): } 35Q74, 35Q86, 35R09.
\end{abstract}




\renewcommand{\theequation}{\thesection.\arabic{equation}}

\section{Introduction}\label{introduction}
\label{sec1}
\setcounter{equation}{0}
This paper aims to study clusters of eigenvalues referred to as the C-ev's of the free oscillation of the Earth (abbreviated by FOE), which we model as a homogeneous isotropic extended Burgers model (abbreviated by EBM). In \cite{YP}, the importance of the EBM emerged in geophysics. The authors of this paper argue that this model can capture the full range of observed mechanical behavior of upper mantle material and reconcile it with the behavior of the Earth's mantle as a whole. For other applications of EBM in geophysics, see \cite{Ivins, LauFaul_2019, Harry} and the references therein. As the first attempt, we provide an analytic method to generate the C-ev's and analysis examining the structure of C-ev's.  

To begin with, we first give a bird's eye view of the C-ev's of FOE before getting into details. The equation of motion for the EBM yields the Boltzmann-type viscoelastic system (abbreviated by BVS). This is an integro-differential system given as the acceleration term 
equal to the divergence of stress over each time interval $[0,t]$ for $t\ge0$. Here, the stress is given as a convolution with a relaxation tensor, which is convoluted with the time derivative of strain. Assuming the strain is zero at the initial time, $t=0$, and integrating the BVS by parts with respect to time, it becomes that the acceleration term is equal to the sum of the elastic term and the memory term. The former and latter terms are referred to as the instantaneous term and memory term, respectively. For simplicity of notation, we assume the density of the EBM is equal to 1 throughout this paper.

At this stage, we replace the two times derivative with respect to time in the BVS by $z^2$-multiplication with $z\in{\mathbb C}$ combined with the traction-free boundary condition on the boundary $\partial\Omega$ of the Earth $\Omega$, assuming it as an open ball in ${\mathbb R}^3$ with radius $R>0$ centered at the origin for simplicity. Then, we have an eigenvalue problem with a parameter $t\ge0$. In this problem, we refer to the integro-differential operator that does not contain the parameter $z$ as a quasi-static EBM (abbreviated by q-EBM). To describe the C-ev's, we first consider the eigenvalues of the boundary value problem for the mentioned instantaneous term in $\Omega$ with the traction-free boundary condition on $\partial\Omega$. Then, for the eigenvectors associated with these eigenvalues, the C-ev's look for the eigenvalues for the mentioned eigenvalue problem for the q-EBM. For further arguments, it is important to note that the elastic tensors of springs in the EBM are commutative because they are isotropic. Based on this, we will show later that these eigenvectors are also the eigenvectors for the eigenvalue problem for the q-EBM with eigenvalues which are the roots of a polynomial of $z$. This is the advantage of capturing the C-ev's because they are roots of a polynomial and we can analyze their structure. 

For simplicity of notation, we assume that the density of $\Omega$ equals $1$ through out the rest of this paper. Before closing this section, we describe the EBM and BVS in details. Figure given after References schematically describes the EBM. Although the EBM is a one space dimensional model, it can be formally generalized to higher space dimensions. The mathematical meaning of the generalization can be given using the spectral decompositions of the elasticity tensors of the springs in the models (see \cite{HKLNT_IPMS, KLLNY} and the references there in). Also, for the isotropic case, there is a mechanical meaning of generalization given in \cite{Flugge} for $n=1$, which is for the Burgers model and yet for the EBM.

To proceed further, we denote a displacement vector for a small deformation of $\Omega$ by $u=u(x,t)\in\R^3$ for a position a.e. $x\in \Omega$ at a time $t\in [0,T)$, where $T\in (0,\infty ]$. We also define the infinitesimal strain tensor by $e[u]:=\frac{1}{2} (\nabla u +(\nabla u)^\mathfrak{t})$
$:~\Omega\times [0,T)\to \R^{3\times 3}_\text{sym}$ with the set ${\mathbb R}^{3\times3}_{\text{sym}}$ of all symmetric matrices and the notation $\mathfrak{t}$ for transposing matrices. Also, the following symbols will be used below.
For matrices $A=(a_{pq}),~B=(b_{pq})\in \R^{3\times 3}$, their inner product is defined by
$A:B:=a_{pq}b_{pq}$, where Einstein's summation symbol is used. For a second-order tensor $\alpha=(\alpha_{pq})\in \R^{3\times 3}$ and
a fourth-order tensor
$C=(c_{pqrs})\in \R^{3\times 3\times 3\times 3}$,
their product is denoted by
$C\alpha:=(c_{pqrs}\alpha_{rs})\in\R^{3\times 3}$.
The product of two fourth order tensors $C=(c_{pqrs})$ and $D=(d_{pqrs})$ is defined by $CD=(c_{pqkl}d_{klrs})$ as well.

Now for the EBM model in Figure, we denote its pair of strain and stress as $(e_i,\sigma_i)\in \R^{3\times 3}_{\text{sym}}\times \R^{3\times 3}_{\text{sym}}$ in accordance with the labeling number $i=0,1,\cdots,n$. Also, we denote by $(e,\sigma)\in \R^{3\times 3}_{\text{sym}}\times \R^{3\times 3}_{\text{sym}}$, the pair of strain and stress of the EBM. Then, using Hooke's law for springs and the fact that the stress of each dashpot is proportional to the instantaneous speed of its strain, the equation of motion for the EBM satisfying the traction-free boundary condition over $\partial\Omega$ is given as follows (see \cite{HKLNT} for the details). 
\begin{equation}\label{s1.1}
\left\{
\begin{array}{ll}
\ddot{u}=\dvg\, \sigma &\mbox{\rm in}~\Omega\times {\mathbb R},\\
\eta_0\dot{\phi}_0=\sigma &    \mbox{\rm in}~\Omega\times {\mathbb R},\\
\eta_i\dot{\phi}_i=\sigma -C_i\phi_i, \quad i=1,\cdots,n &    \mbox{\rm in}~\Omega\times {\mathbb R},\\
\sigma\nu=0 &\mbox{\rm on}~\partial\Omega\times {\mathbb R},\\
%(u, \dot{u}, \psi, \phi)|_{t=0}=(u^0, \dot{u}^0,0, 0)& \mbox{\rm in}~\Omega
\end{array}
\right.
\end{equation}
with the divergence operator $\text{div}$, the unit outernormal of $\partial\Omega$, and the total stress tensor $\sigma$ defined by
\begin{align}\label{s1.2}
\sigma :=C_0\psi,\quad
  \psi:=
  e[u]-\sum_{i=0}^n\phi_i=e_0^s,
\end{align}
where for each $i\,(0\le i\le n)$ corresponding to the label $i$ of Figure. Also, $\eta_i>0$ and $C_i=((C_i)_{pqrs}):=({\lambda}_i\delta_{pq}\delta_{rs}+{\mu}_i(\delta_{pr}
\delta_{qs}+\delta_{ps}\delta_{qr}))$ with Kronecker's delta $\delta_{ij}$ are the inhomogeneous viscosity and elasticity tensor of the dashpot and spring, respectively. Whenever we use the terminology inhomogeneous in this paper, the object depending on the space variable $x\in\Omega$ is bounded measurable in a bounded Lipschitz domain $\Omega\subset{\mathbb{R}}^3$, and for notational simplicity, we suppress $x$. For the FOE, $\Omega$ is the mentioned ball, and the physical objects there are all homogeneous. Also, for each $\phi_i$ with $i=0,\cdots,n$ is the strain of the $i$-th dashpot, and we denote $\phi:=(\phi_1,\cdots,\phi_n)$, 
and the time derivative by $\dot{}$ for notational convenience. Here, we assume that the Lam\'e moduli $\lambda_i,\,\mu_i$ satisfy the strong convexity condition:
\begin{equation}\label{strong convex}
\mu_i\ge\delta,\,\,3\lambda_i+2\mu_i\ge\delta,\,\,i=0,\cdots,n
\end{equation}
for some constant $\delta>0$.
Next, by eliminating $(\phi_1,\cdots,\phi_n)$ using the second, third equations of \eqref{s1.1} and assuming $(\phi_1,\cdots,\phi_n)=0$ at $t=0$, we uniquely derive a relaxation tensor $G(t)$ satisfying the causality, $G(t)=0$ for $t<0$, also it is completely monotone fourth order tensor acting on $L^\infty(\Omega;{\mathbb R}^{3\times3}_{\text{sym}})$, the set of ${\mathbb R}^{3\times3}_{\text{sym}}$ valued bounded measurable functions in $\Omega$, over the time interval $[0,\infty)$ satisfying $G(0)=C_0$, and the stress-strain relation can be given as
\begin{equation}\label{stress_strain rel}
\sigma(t,x)=\int_0^t G(t-s)e[\dot{u}](s,x)\,ds
\end{equation}
(see \cite{HKLNT, Zemanian} for the details by viewing the EBM as a one-port mechanical system). Here, note that from the causality of $G(t)$ and \eqref{stress_strain rel}, we could have assumed $\psi=0$ at $t=0$. Henceforth, we will always take this into account, so that we assume
\begin{equation}\label{zero initial strains}
\phi:=(\phi_1,\cdots,\phi_n)=0,\,\,\psi=0\,\,\,\text{at}\,\,\,t=0.
\end{equation}

Then, combining this with the first equation of \eqref{s1.1}, we have
the BVS as follows:
\begin{equation}\label{Boltzmann_1}
\ddot{u}(t,x)=\dvg\int_0^t G(t-s)e[\dot{u}](s,x)\,ds.   \end{equation}
Further, assuming $e[u](t,x)=0$ at $t=0$, we have
\begin{equation}\label{Boltzmann_2}
\ddot{u}=\dvg\{C_0 e[u](t,x)\}+\int_0^t \dot{G}(t-s)e[u](s,x)\,ds.
\end{equation}
with the instantaneous term $\dvg\{C_0 e[u](t,x)\}$ and the q-EBM operator $\dvg\{C_0 e[u](t,x)\}+\int_0^t \dot{G}(t-s)e[u](s,x)\,ds$ acting to $u(t,x)$.

\medskip
The rest of this paper is organized as follows.
In the next section, we focus on the method giving an explicit form of the relaxation kernel of the inhomogeneous isotropic EBM defined in $\Omega\subset{\mathbb R}^3$. Hence, in these two sections, the elasticity tensors and viscosities can depend on $x\in\Omega$, and they are bounded measurable in $\Omega$. Section 3 is devoted to finding some special radial eigenfunctions for the instantaneous term which are also eigenfunctions of the q-EBM. In Section 4, we provide a method to generate C-ev's for the FOE. Then in Section 5, we analytically analyze the structure of C-ev's for the FOE and give an inversion formula identifying the q-EBM by knowing the C-ev's. The final section, Section 6, is for the conclusions and discussions.






\begin{comment}
Let $\Omega$ be a bounded domain in a $d$-dimensional Euclidean space ${\mathbb R}^d$ with $d=2,3$ and Lipschitz smooth boundary $\Gamma:=\partial\Omega$. We denote the outward unit normal of $\Gamma$
by $\nu\in \R^d$. As we mentioned before, we consider $\Omega$ as a reference domain on which we consider a small viscoelastic deformation. We denote a displacement for a small deformation of $\Omega$ by $u=u(x,t)\in\R^d$ for a position a.e. $x\in \Omega$ and a time $t\in [0,T)$, where $T\in (0,\infty ]$. We also define the infinitesimal strain tensor by $e[u]:=\frac{1}{2} (\nabla u +\nabla u^T)$
$:~\Omega\times [0,T)\to \R^{d\times d}_\text{sym}$.

The following symbols will be used below.
For matrices $A=(a_{pq}),~B=(b_{pq})\in \R^{d\times d}$, their inner product is defined by
$A:B:=a_{pq}b_{pq}$, where Einstein's summation symbol is used. For a matrix $\xi=(\xi_{pq})\in \R^{d\times d}$ and
a fourth-order tensor
$C=(c_{pqrs})\in \R^{d\times d\times d\times d}$,
their product is denoted by
$C\xi:=(c_{pqrs}\xi_{rs})\in\R^{d\times d}$.
The product of two fourth order tensors $C=(c_{pqrs})$ and $D=(d_{pqrs})$ is defined by $CD=(c_{pqkl}d_{klrs})$ as well.
Also, the set of $d$-dimensional symmetric matrices is represented by $\R^{d\times d}_\text{sym}$.
For a fourth-order tensor
$C=(c_{pqrs})$, it is said that the major symmetry condition is satisfied when $c_{pqrs}=c_{rspq}$ holds and that the minor symmetry condition is satisfied when $c_{pqrs}=c_{qprs}=c_{pqsr}$
holds, for all $p,q,r,s=1,\cdots,d$.
When the major and minor symmetries are both satisfied,
it is called full symmetry.
We remark that a linear map from $\R^{d\times d}_\text{sym}$ into itself is uniquely represented by
a fourth-order tensor $C$ with the minor symmetry
as $\R^{d\times d}_\text{sym}\ni \xi\mapsto C\xi\in \R^{d\times d}_\text{sym}$.


Corresponding to each of these models, we denote its pair of strain and stress as $(e_i,\sigma_i)\in \R^{d\times d}_{\text{sym}}\times \R^{d\times d}_{\text{sym}}$ in accordance with the labeling number $i=0,1,\cdots,n$. Also, we denote by $(e,\sigma)\in \R^{d\times d}_{\text{sym}}\times \R^{d\times d}_{\text{sym}}$, the pair of strain and stress of the EBM. Then, assuming that the strain and stress tensors depend on a.e. $x\in \Omega$ and $t\in [0,T)$, their relation for each model and the equation of motion for the EBM are given as follows.



\begin{align}
M_0~&:~
\begin{cases}
  \sigma_0^s=C_0e_0^s,\quad \eta_0\dot{e}_0^d=\sigma_0^d\\
  e_0=e_0^s+e_0^d,\quad \sigma_0=\sigma_0^s=\sigma_0^d
\end{cases} \label{M0}\\
K_i~&:~
\begin{cases}
  \sigma_i^s=C_ie_i^s,\quad \eta_i\dot{e}_i^d=\sigma_i^d\\
  e_i=e_i^s=e_i^d,\quad \sigma_i=\sigma_i^s+\sigma_i^d
\end{cases}
\quad (i=1,\cdots,n)\label{Ki}\\
\mbox{EBM}~&:~
\begin{cases}
  e=e[u]\\
  \rho \ddot{u}=\dvg \sigma \\
  \DS{e=\sum_{i=0}^ne_i,\quad \sigma=\sigma_0=\sigma_1=\cdots =\sigma_n}.
\end{cases}\label{EBM}
\end{align}
Here, we recall that "$\cdot$" stands for the $t$-derivative.

For each $i=0,1,\cdots,n$, $(e_i^s, \sigma_i^s)$ and $ (e_i^d, \sigma_i^d)$ are the pairs of strain and stress for the spring and dashpot, respectively. Also, $C_i\in \R^{d\times d\times d\times d}$ and $\eta_i>0$ for each $i=0,1,\cdots,n$ are the elasticity tensor of the spring and viscosity of the dashpot, respectively.


\medskip
We assume the following conditions for $\eta_i$'s and $C_i$'s.


\begin{assumption}${}$
\begin{itemize}
\item [{\rm(i)}] $\eta_i\in L^\infty(\Omega;\R),\,\,C_i\in L^\infty(\Omega; \R^{d\times d\times d\times d}),\,i=0,1,\cdots,n$. That is, for each $i=0,1,\cdots,n\,$, $\eta_i$ and $C_i$ are $\R$-valued and $\R^{d\times d\times d\times d}$-valued bounded measurable function defined in $\Omega$, respectively.
\item [{\rm(ii)}] There exists $\eta_{\ast}>0$ such that each $\eta_i$ satisfies $\eta_i(x)\ge \eta_{\ast},\, \text{a.e.}\, x\in\Omega$.
\item [{\rm(iii)}] Each $C_i(x)=(c^{(i)}_{pqrs}(x))$ satisfies the full symmetry condition:
$$c^{(i)}_{pqrs}(x)=c^{(i)}_{rspq}(x)=c^{(i)}_{qprs}(x),\,\,\text{a.e.}\,x\in\Omega.$$
\item [{\rm (iv)}] Each $C_i(x)=(c^{(i)}_{pqrs}(x))$ satisfies the strong convexity condition:
\[
(C_i(x)\xi):\xi \ge c_*|\xi|^2,\,\,\xi\in\R^{d\times d},\,\, \text{a.e.}\,x\in\Omega,~,
\]
for some $c_*>0$. where
\[
\sum_{p,q,r,s=1}^d c^{(i)}_{pqrs}(x)\xi_{pq}\xi_{rs}\ge c_*|\xi|^2, \,\,\xi\in\R^{d\times d},\,\, \text{a.e.}\,x\in\Omega.
\]
\end{itemize}
\end{assumption}

\medskip
We set $\phi_i:=e_i^d$ for $i=0,\cdots,n$
and define $\phi:=(\phi_1,\cdots,\phi_n)$
$:~\Omega\times [0,T)\to (\R^{d\times d}_\text{sym})^n$, $\bar\phi:=(\phi_0,\phi)=(\phi_0,\cdots,\phi_n)$
$:~\Omega\times [0,T)\to (\R^{d\times d}_\text{sym})^{n+1}$
and also set $\psi :=e_0^s$
$:~\Omega\times [0,T)\to \R^{d\times d}_\text{sym}$.
It should be remarked here that throughout this paper, spatial variables will not be specified unless necessary. Also, $x\in\Omega$ or $\text{a.e.}\,x\in\Omega$ will not be mentioned as long as it can be determined from the context.

Then, the conditions \eqref{M0}, \eqref{Ki}, \eqref{EBM} are equivalent to
\begin{align*}
M_0~&:~
  \sigma=C_0\psi,\quad \eta_0\dot{\phi}_0=\sigma \\
K_i~&:~
\eta_i\dot{\phi}_i=\sigma -C_i\phi_i \quad (i=1,\cdots,n)\\
\mbox{EBM}~&:~
\begin{cases}
  \DS{e[u]=\psi+\sum_{i=0}^n\phi_i}\\
  \rho \ddot{u}=\dvg \sigma
\end{cases}
\end{align*}

Now, let $\Gamma_N$ be a non-empty open subset of $\partial \Omega$ and define $\Gamma_D:=\partial\Omega\setminus\overline{\Gamma_N}$. If $d=3$ and $\Gamma_D\not=\emptyset$, we assume that $\overline{\Gamma_D}\cap\overline{\Gamma_N}$ is Lipschitz smooth. Here, for a set $S\subset\R^d$,  we denoted its closure by $\overline S$.
Assume that the displacement $u$ is fixed at $\Gamma_D$, the traction at $\Gamma_N$ is free, and no external force is applied to $\Omega$.

Then, the initial boundary value problem for the EBM with an initial data $(u^0,\dot u^0,\bar\phi^0)$ for $(u,\dot u,\bar\phi)$ at the initial time $t=0$ is to find $(u.\bar\phi):\Omega\times [0,T)\rightarrow \R^d\times(\R^{d\times d}_\text{sym})^{n+1}$ which satisfies

\begin{equation}\label{B2.4}
\left\{
\begin{array}{ll}
\rho \ddot{u}=\dvg\, \sigma &\mbox{\rm in}~\Omega\times (0,T),\\
\eta_0\dot{\phi}_0=\sigma &    \mbox{\rm in}~\Omega\times (0,T),\\
\eta_i\dot{\phi}_i=\sigma -C_i\phi_i \quad (i=1,\cdots,n)&    \mbox{\rm in}~\Omega\times (0,T),\\
u=0&\mbox{\rm on}~\Gamma_D\times (0,T),\\
\sigma\nu=0 &\mbox{\rm on}~\Gamma_N\times (0,T),\\
(u, \dot{u}, \bar\phi)|_{t=0}=(u^0, \dot{u}^0, \bar\phi^0)& \mbox{\rm in}~\Omega,
\end{array}
\right.
\end{equation}
where the total stress tensor $\sigma$ is defined by
\begin{align}\label{B2.5}
\sigma :=C_0\psi,\quad
  \psi:=
  e[u]-\sum_{i=0}^n\phi_i.
\end{align}

The importance of the relaxation tensor is that it converts the EBM to the BVS. To begin with, we recall that the relaxation tensor $G(x,t,\tau)\in \R^{d\times d\times d\times d}$
for $\text{a.e. $x\in\Omega$}$ and $0\le s\le t<T$ is a fourth-order tensor, and that it is the kernel that gives the constitutive equation
\begin{equation}\label{B2.6}
\sigma(x,t)=\int_0^t G(x,t,s)\dot e(x,s)\,ds,
\end{equation}
where $\sigma(x,t)$ and $e(x,t)$
represent the stress tensor and strain tensor, respectively. Note that
\eqref{B2.6} holds on $E=\{(x,t,s):\,0\le s\le t,\,\,\text{a.e $x\in\Omega$}\}$ for any $e(x,s)$. Due to the symmetry of the strain tensor $e(x,s)$ and stress tensor $\sigma(t,s)$, it is natural to assume that $G(x,t,s)$ satisfies the minor symmetry. We further assume that $G(x,t,s)$ is continuous with respect to $(t,s)$ for each fixed $\text{a.e.}\,x\in\Omega$. Summarizing what has been mentioned above, we make the following assumptions.

\begin{assumption}\label{assumption for G}${}$
\begin{itemize}
\item [{\rm (i)}]
$G(x,t,s)$ is defined on $E=\{(x,t,s):\,0\le s\le t,\,\text{a.e. $x\in\Omega$}\}$ and it is continuous with respect to $(t,s)$ fore each fixed $\text{a.e. $x\in\Omega$}$.
\item [{\rm (ii)}]
The foruth-order tensor $G(x,t,s)$ satisfies the minor symmetry.
\end{itemize}
\end{assumption}

\bigskip
Then, we have the following uniqueness for the relaxation tensor.

\begin{Th}\label{uniqueness of relaxation tensor}
Under Assumption \ref{assumption for G}, there is at most one relaxation tensor.
\end{Th}
\pf $\,{}$
Let $G(x,t,s),\,G'(x,t,s)$ satisfy \eqref{B2.6}.
Then, we have
\begin{equation}\label{B2.7}
\int_0^t (G(x,t,s)-G'(x,t,s))\dot e(x,s)\,ds=0\,\,\,\text{a.e. $x\in\Omega$ for any
strain $e(x,s)$.}
\end{equation}
Suppose that there exists $(x_0,t_0,s_0)\in E$ such that $G(x_0,t_0,s_0)\not=G'(x_0,t_0,s_0)$. Here, by the continuity of $G(x,t,s),\,G'(x,t,s)$ with respect to $(t,s)$, we can assume $0<s_0\le t_0$. Then, there exist $i,j,k,l\in\{1,2,\cdots,d\}$ such that $G_{ijkl}(x_0,t_0,s_0)\not=G'_{ijkl}(x_0,t_0,s_0)$,
From \eqref{B2.7}, we have
\begin{equation}\label{B2.8}
\int_0^t \displaystyle\sum_{p,q=1}^d (G_{ijpq}(x,t,s)-{G'}_{ijpq}(x,t,s))\dot e_{pq}(x,s)\,ds=0\,\,\,\text{a.e. $x\in\Omega$ for any
 $(e_{pq}(x,s))$.}
\end{equation}

\bigskip
Since $e=e[u]$ and $\dot{e}=e[\dot{u}]$, an initial and boundary value problem
of the BVS with the constitutive equation \eqref{B2.6} of the convolutional type is given as follows:
\begin{equation}\label{B2.9}
\left\{
\begin{array}{ll}
\rho \ddot{u}=\dvg\, \sigma &\mbox{\rm in}~\Omega\times (0,T),\\
{\DS \sigma(x,t)=\int_0^t G(x,t-s)e[\dot{u}](x,s)\,ds} & \mbox{\rm in}~\Omega\times (0,T),\\
u=0&\mbox{\rm on}~\Gamma_D\times (0,T),\\[5pt]
\sigma\nu=0 &\mbox{\rm on}~\Gamma_N\times (0,T),\\[5pt]
(u, \dot{u})|_{t=0}=(u^0, \dot{u}^0)& \mbox{\rm in}~\Omega.
\end{array}
\right.
\end{equation}

 We remark here that if the relaxation is of the convolution type and has the exponentially decaying property  with some uniform estimate, then solutions of the initial boundary value problem \eqref{B2.9} with the mentioned constrained on $G$ and have the same property.
From the second equation of \eqref{B2.9}, we have
\begin{equation}\label{B2.10}
\begin{aligned}
 \sigma(x,t)&=\int_0^t G(x,t-s)e[\dot{u}](x,s)\,ds\\
 &=G(x,0)e[u](x,t)+\int_0^t \dot{G}(x,t-s)e[u](x,s)ds
\end{aligned}
\end{equation}
with $\dot G(x,t):=\partial_t G(x,t)$ because it is usually assumed that $e[u]=0$ at $t=0$.
Here, $G(x,0)$ is called the instantaneous elastic tensor.

\bigskip
To transform the EBM \eqref{B2.4} to the form of BVS \eqref{B2.9},
we rewrite \eqref{B2.4} by using $(u,\psi,\phi)$ instead of $(u,\bar\phi)$.
From \eqref{B2.4} and \eqref{B2.5}, we obtain
\begin{align*}
\dot\psi&=e[\dot{u}]-\sum_{i=0}^n\dot\phi_i
=e[\dot{u}]-\eta_0^{-1}\sigma -\sum_{i=1}^n\eta_i^{-1}(\sigma -C_i\phi_i)\\
&=e[\dot{u}]-(\sum_{i=0}^n\eta_i^{-1})\sigma +\sum_{i=1}^n\eta_i^{-1}C_i\phi_i
=e[\dot{u}]-(\sum_{i=0}^n\eta_i^{-1})C_0\psi+\sum_{i=1}^n\eta_i^{-1}C_i\phi_i.
\end{align*}

Combining this with \eqref{B2.4}, we have the following system for $(\psi,\phi)$
\begin{equation}\label{B2.11}
\left\{
\begin{aligned}
%\partial_t u&=v\\
%\partial_t v&=\rho^{-1}\dvg C_0\psi\\
\partial_t\psi&=e[\dot{u}]-(\sum_{i=0}^n\eta_i^{-1})C_0\psi+\sum_{i=1}^n\eta_i^{-1}C_i\phi_i\\
\partial_t\phi_i&=\eta_i^{-1}C_0\psi-\eta_i^{-1}C_i\phi_i, \quad 1\leq i\leq n.
\end{aligned}
\right.
\end{equation}
This yields

\begin{equation}\label{B2.12}
\begin{pmatrix}
\psi \\
\phi
\end{pmatrix}
(x,t)=e^{tL(x)}
\begin{pmatrix}
\psi^0 \\
\phi^0
\end{pmatrix}
+\int_0^t e^{(t-s)L_b(x)}
\begin{pmatrix}
e[\dot{u}] \\
0
\end{pmatrix}
(x,s)\, ds,
\end{equation}
where
\begin{equation}\label{B2.13}
L(x):=
\begin{pmatrix}
                       -(\sum_{i=0}^n\eta_i^{-1})C_0
                       & \eta_1^{-1}C_1 & \cdots     & \eta_{n-1}^{-1}C_{n-1}            &    \eta_n^{-1}C_n    \\

                             \eta_1^{-1}C_0                              & -\eta_1^{-1}C_1 & \cdots     &    0               &            0         \\
            \cdots                                   & \cdots                   & \cdots     & \cdots             & \cdots               \\
                               \eta_{n-1}^{-1}C_0                               & 0              & \cdots     & -\eta_{n-1}^{-1}C_{n-1}     &            0         \\
                              \eta_n^{-1}C_0                              & 0              & \cdots     &    0               &    -\eta_n^{-1}C_n
\end{pmatrix},
\end{equation}
and $L(x)$ is a linear map from $(\R^{d\times d}_\text{sym})^{n+1}$ into itself for a fixed $x\in\Omega$.
We consider the "block matrix" representation of the operator
$e^{tL_b(x)}$ as
\begin{equation}\label{B2.14}
e^{tL_b(x)}=
\begin{pmatrix}
G_{00}(x,t)&G_{01}(x,t)&\cdots&G_{0n}(x,t)\\
\ldots& & &\ldots\\
G_{n0}(x,t)&G_{n1}(x,t)&\cdots&G_{nn}(x,t)
\end{pmatrix},
\end{equation}
where each "matrix" element $G_{ij}(x,t)$ is a
fourth-order tensor with minor symmetry.
From \eqref{B2.12}, we obtain
\begin{align*}
&\psi(x,t)=\int_0^tG_{00}(x,t-s)e[\dot{u}](x,s)\,ds +\psi_*(x,t),\\
&\psi_*(x,t):=G_{00}(x,t)\psi^0(x)+
\sum_{i=1}^nG_{0i}(x,t)\phi_i^0(x).
\end{align*}
Since
\[
\psi^0=\psi(\cdot,0)=e[u^0]-\sum_{i=0}^n\phi_i^0
=e[u^0]-\phi_0^0-\sum_{i=1}^n\phi_i^0,
\]
if $\phi_0^0=e[u^0]$ and $\phi_1^0=\cdots =\phi_n^0=0$,
then $\psi_*=0$ holds.
Assuming that $\psi_*=0$ holds, from $\sigma=C_0\psi$, we obtain the
constitutive relation of convolutional Boltzman-type \eqref{B2.6} where
\begin{equation}\label{B2.15}
G(x,t-s):=C_0G_{00}(x,t-s)\,\,\text{on}\,\,E=\{(x,t,s):\,0\le s\le t,\,\,\text{a.e. $x\in\Omega$}\}.
\end{equation}
This $G$ is a convolutional relaxation tensor with minor symmetry
and \eqref{B2.4} is equivalent to \eqref{B2.9}. Combining \eqref{B2.6}, \eqref{B2.10} and \eqref{B2.15}, we have that
\begin{equation}\label{B2.16}
\begin{aligned}
 \sigma(x,t)&=\int_0^t G(x,t-s)e[\dot{u}](x,s)\,ds\\
 &=C_0e[u](x,t)+\int_0^t C_0\dot{G}_{00}(x,t-s)e[u](x,s)ds.
\end{aligned}
\end{equation}

\section{The relaxation tensor-r}\label{relaxtion tensor-r}\setcounter{equation}{0}

In this section, we will focus on computing $C_0\dot{G}_{00}(x,t-s)e[u](x,s)$. Since $G_{00}(x,t)$ is the $(1,1)$ component of $e^{tL(x)}=\sum_{k=0}t^k\frac{L^k(x)}{k!}$, we first compute $L^k(x)$. Denote the $(1,1)$ component of $L^k(x)$ by $L_{00k}(x)$. 
We divide $L$ into $L=\breve{D}+\breve{M}+\breve{N}$,
where 
\begin{equation}\label{r7.1}
\breve{D}:=
\begin{pmatrix}
                       -(\sum_{i=0}^n\eta_i^{-1})C_0
                       & 0 & \cdots     & 0            &    0    \\

                             0                              & -\eta_1^{-1}C_1 & \cdots     &    0               &            0         \\
            \cdots                                   & \cdots                   & \cdots     & \cdots             & \cdots               \\
                               0                               & 0              & \cdots     & -\eta_{n-1}^{-1}C_{n-1}     &            0         \\
                              0                              & 0              & \cdots     &    0               &    -\eta_n^{-1}C_n
\end{pmatrix},
\end{equation}
\begin{equation}\label{r7.2}
\breve{M}:=
\begin{pmatrix}
                      0
                       & \eta_1^{-1}C_1 & \cdots     & \eta_{n-1}^{-1}C_{n-1}            &    \eta_n^{-1}C_n    \\

                             0                            &0 & \cdots     &    0               &            0         \\
            \cdots                                   & \cdots                   & \cdots     & \cdots             & \cdots               \\
                               0                               & 0              & \cdots     & 0     &            0         \\
                              0                            & 0              & \cdots     &    0               &    0
\end{pmatrix},
\end{equation}

\begin{equation}\label{r7.3}
\breve{N}:=
\begin{pmatrix}
                       0
                       & 0 & \cdots     & 0            &    0    \\

                             \eta_1^{-1}C_0                              & 0     &    0   &      0    & 0     \\
            \cdots                                   & \cdots                   & \cdots     & \cdots             & \cdots               \\
                               \eta_{n-1}^{-1}C_0                               & 0              & \cdots     &0    &            0         \\
                              \eta_n^{-1}C_0                              & 0              & \cdots     &    0               &    0
\end{pmatrix}.
\end{equation}
Define the sets $\mathcal{D}$, $\mathcal{M}$  $\mathcal{N}$ and $\mathcal{R}$ by
\begin{equation}\label{r7.4}
\mathcal{D}:=
\begin{pmatrix}
* & 0 & \cdots     & 0            &    0    \\

0     & *     &    0      &    0       &          0       \\
            \cdots                                   & \cdots                   & \cdots     & \cdots             & \cdots               \\
                               0                               & 0              & \cdots     & *     &            0         \\
                              0                              & 0              & \cdots     &    0               &    *
\end{pmatrix},
\end{equation}
\begin{equation}\label{r7.5}
\mathcal{M}:=
\begin{pmatrix}
 0  & *     & *      & *      &    *   \\

0   &   0   & \cdots     &    0      &      0      \\
            \cdots                                   & \cdots                   & \cdots     & \cdots             & \cdots               \\
                               0                               & 0              & \cdots     & 0     &            0         \\
                              0                            & 0              & \cdots     &    0               &    0
\end{pmatrix},
\end{equation}

\begin{equation}\label{r7.6}
\mathcal{N}:=
\begin{pmatrix}
0  & 0 & \cdots     & 0       &    0    \\

*  & 0     &    0       &  0     &      0     \\
\cdots      & \cdots  & \cdots     & \cdots     & \cdots   \\
*      & 0    & \cdots   &  0    &       0       \\
*      & 0    & \cdots     &    0        &    0
\end{pmatrix},
\end{equation}


\begin{equation}\label{r7.7}
\mathcal{R}:=
\begin{pmatrix}
0  & 0 & \cdots     & 0       &    0    \\

0  & *     &    *       &  *     &      *     \\
0     & *  & \cdots     & \cdots     & \cdots   \\
0     & *   & \cdots   &  *   &       *       \\
0     & *    & \cdots   & *   &    *
\end{pmatrix}.
\end{equation}

We note that $\breve{D}\in \mathcal{D}$, $\breve{M}\in \mathcal{M}$ and $\breve{N}\in \mathcal{N}$.
We divide $\mathcal{D}$, $\mathcal{M}$, $\mathcal{R}$ and $\mathcal{N}$ into two groups by its $(1,1)$ block element. Let us define $\mathcal{O}$ to be the set having  $0$ in $(1,1)$ block element and $\mathcal{O}^c$ to be the set having  nonzero $(1,1)$ block element. Thus, we know that $M,N,R\subset \mathcal{O}$.

\begin{lemma}\label{B_lem7.1}
We have $MD\in M$, $DM\in M$, $ND\in N$, $DN\in N$, $RD\in R$, $DR\in R$ and $DD\in D$. The product with $\mathcal{D}$ preserves the original type.
\end{lemma}

\begin{lemma}\label{B_lem7.2}
We have $NN=0$, $NM\in R$, $MM=0$, $MN\in D$, $RN\in N$, $RM=0$ and $RR\in R$. Here, only $MN\in \mathcal{O}^c$.
\end{lemma}
There are $3^k$ terms in $L^k=(\breve{D}+\breve{M}+\breve{N})^k$. From 
Lemma \ref{B_lem7.1} and Lemma \ref{B_lem7.2}, only the terms belongs to $\mathcal{D}$ will be in $\mathcal{O}^c$. Thus, we define $\mathcal{D}_{2q}^k$ to be the terms in $L^k=(\breve{D}+\breve{M}+\breve{N})^k$ belongs to $\mathcal{D}$ having exactly $q$ elements of $\breve{M}$ and $\breve{N}$.

\begin{lemma}\label{B_lem7.3}
$L_{00k}(x)=\sum_{q=0}^{[k/2]}\mathcal{D}_{2q}^k$.
\end{lemma}

\begin{Th}\label{B_thm7.4}
Assume that  $\breve{D}\breve{M}=\breve{M}\breve{D}$; that is $bC_0=\eta_j^{-1}C_j$ for $1\leq j\leq n$. Then
\begin{equation}\label{B7.8}
\begin{aligned}
G_{00}(x,t)=\frac{1}{2}\{e^{-t(bC_0+\sqrt{\breve{M}\breve{N}}\,)}+e^{-t(bC_0-\sqrt{\breve{M}\breve{N}}\,)}\},
\end{aligned}
\end{equation}
where $b=\sum_{i=0}^n\eta_i^{-1}$.
\end{Th}
\begin{proof}
From Lemma \ref{B_lem7.3}, we have that
\begin{equation}\label{B7.9}
\begin{aligned}
G_{00}(x,t)=&\sum_{k=0}^{\infty}\frac{1}{k!}t^kL_{00k}(x)\\
=&\sum_{k=0}^{\infty}\frac{1}{k!}t^k\sum_{q=0}^{[k/2]}\mathcal{D}_{2q}^k\\
=&\sum_{k=0}^{\infty}\frac{1}{k!}t^k\sum_{q=0}^{[k/2]}\frac{k!}{(2q)!(k-2q)!}(-bC_0)^{k-2q}\cdot(\breve{M}\breve{N})^q\\
=&\sum_{k=0}^{\infty}\frac{1}{k!}(-t)^k\sum_{q=0}^{[k/2]}\frac{k!}{(2q)!(k-2q)!}(bC_0)^{k-2q}\cdot (\sqrt{\breve{M}\breve{N}})^{2q}.
\end{aligned}
\end{equation}
By using the binomial theorem, we have that
\begin{equation}\label{B7.10}
\begin{aligned}
2\sum_{q=0}^{[k/2]}\frac{k!}{(2q)!(k-2q)!}(bC_0)^{k-2q}\cdot (\sqrt{\breve{M}\breve{N}})^{2q}=(bC_0+\sqrt{\breve{M}\breve{N}}\,)^k+(bC_0-\sqrt{\breve{M}\breve{N}})^k.
\end{aligned}
\end{equation}
Combining \eqref{B7.9} and \eqref{B7.10}, we get that
\begin{equation}\label{B7.11}
\begin{aligned}
G_{00}(x,t)=&\sum_{k=0}^{\infty}\frac{1}{k!}(-t)^k\sum_{q=0}^{[k/2]}\frac{k!}{(2q)!(k-2q)!}(bC_0)^{k-2q}\cdot (\sqrt{\breve{M}\breve{N}})^{2q}\\
=&\frac{1}{2}\sum_{k=0}^{\infty}\frac{1}{k!}(-t)^k\{(bC_0+\sqrt{\breve{M}\breve{N}}\,)^k+(bC_0-\sqrt{\breve{M}\breve{N}})^k\}\\
=&\frac{1}{2}\{e^{-t(bC_0+\sqrt{\breve{M}\breve{N}}\,)}+e^{-t(bC_0-\sqrt{\breve{M}\breve{N}}\,)}\}.
\end{aligned}
\end{equation}
\end{proof}

\end{comment}

\section{Derivation of the relaxation tensor for inhomogeneous isotropic q-EBM }\label{formula I-EBM}
\setcounter{equation}{0}
In this section, we derive an explicit form of the relaxation tensor for inhomogeneous isotropic EBM. Note that we already gave a general method to derive the relaxation tensor for inhomogeneous anisotropic EBM in \cite{HKLNT}. However, the form of the relaxation tensor was not explicit enough to study the FOE. We would like to emphasize that this paper gave an explicit form of the relaxation tensor for three-dimensional inhomogeneous isotropic EBM for the first time. For the derivation given below, we use the volumetric/deviatoric decompositions and easily computable symmetric tensor block matrices \eqref{s2.11}/\eqref{s2.13} associated with the ordinary differential systems \eqref{s2.8}/\eqref{s2.9} for the volumetric/deviatoric $\psi$, $\phi_i$'s. We note that without using the decompositions, the derivation quickly hits a snag.

To begin with, we first introduce volumetric projection $I_m$ and deviatoric projection $J_m$ acting on $L^\infty(\Omega;{\mathbb R}^{3\times3}_{\text{sym}})$ as multiplication by 4th order tensors which synchronize well with the EBM under consideration. They are commutative projection operators on $L^\infty(\Omega;{\mathbb R}^{3\times3}_{\text{sym}})$ defined as multiplication by the following fourth-order tensors
\begin{equation}\label{I-J}
I_m:=(3^{-1}\delta_{ij}\delta_{kl}),\,\,J_m:=\tilde I-I_m,    
\end{equation}
where $\tilde I:=I_3\otimes I_3$ is a fourth-order tensor with the $3\times3$ identity matrix $I_3$. Here, the multiplications between $\tilde I$ and the second-order tensors, fourth-order tensors follow those for tensors defined in Section \ref{introduction}. Since $I_m\,J_m=0$, the ranges of $I_m$ and $J_m$ are orthogonal.
Then, for $0\leq j\leq n$, we have
\begin{equation}\label{s2.2}
\left\{
\begin{array}{l}
\phi^V_j:=I_m\, \phi_j=3^{-1}{\rm tr} (\phi_j)\,I_3,\\
\phi^D_j:=J_m\,\phi_j=\phi_j-\phi^V_j
\end{array}
\right.
\end{equation}
and
\begin{equation}\label{s2.3}
\left\{
\begin{array}{l}
\psi^V:=I_m\,\psi==3^{-1}{\rm tr} (\psi)\,I_3,\\
\psi^D:=J_m\,\psi=\psi-\psi^V,
\end{array}
\right.
\end{equation}
where $I_3$ is the $3\times3$ identity matrix, and $\text{tr}(\phi_j)$ denotes the trace of $\phi_j$.

Now, combining \eqref{s1.1} and \eqref{s1.2}, we have
\begin{equation}\label{s2.1}
\left\{
\begin{aligned}
%\partial_t u&=v\\
%\partial_t v&=\rho^{-1}\dvg C_0\psi\\
\partial_t\psi&=e[\dot{u}]-(\sum_{i=0}^n\eta_i^{-1})C_0\psi+\sum_{i=1}^n\eta_i^{-1}C_i\phi_i\\
\partial_t\phi_i&=\eta_i^{-1}C_0\psi-\eta_i^{-1}C_i\phi_i, \quad 1\leq i\leq n.
\end{aligned}
\right.
\end{equation}
Rewrite \eqref{s2.1} as
\begin{equation}\label{s2.4}
\left\{
\begin{aligned}
\partial_t\psi^V+\partial_t\psi^D=&e[\dot{u}]-bC_0\psi+\sum_{i=1}^n\eta_i^{-1}C_i\phi_i\\
=&e[\dot{u}]-b(3\lambda_0+2\mu_0)\psi^V-2b\mu_0\psi^D\\
&+\sum_{i=1}^n\eta_i^{-1}(3\lambda_i+2\mu_i)\phi_i^V+\sum_{i=1}^n2\mu_i\eta_i^{-1}\phi_i^D,\\
\partial_t\phi_i^V+\partial_t\phi_i^D=&\eta_i^{-1}(3\lambda_0+2\mu_0)\psi^V+2\eta_i^{-1}\mu_0\psi^D\\
&-\eta_i^{-1}(3\lambda_i+2\mu_i)\phi_i^V-2\mu_i\eta_i^{-1}\phi_i^D, \quad 1\leq i\leq n.
\end{aligned}
\right.
\end{equation}
Note here that we have
\begin{equation}\label{I_J e}
I_m\,e[\dot u]=3^{-1}(\nabla\cdot \dot u)I_3,\,\, J_m\,e[\dot u]=e[\dot u]-3^{-1}(\nabla\cdot \dot u)I_3.
\end{equation}
\begin{comment}
Taking trace on \eqref{s2.1}, we obtain 
\begin{equation}\label{s2.5}
\left\{
\begin{aligned}
\partial_t({\rm tr}(\psi))=&\nabla\cdot\dot{u}-b(3\lambda_0+2\mu_0){\rm tr}(\psi)+\sum_{i=1}^n\eta_i^{-1}(3\lambda_i+2\mu_i){\rm tr}(\phi_i),\\
\partial_t({\rm tr}(\phi_i))=&\eta_i^{-1}(3\lambda_0+2\mu_0){\rm tr}(\psi)-\eta_i^{-1}(3\lambda_i+2\mu_i){\rm tr}(\phi_i), \quad 1\leq i\leq n.
\end{aligned}
\right.
\end{equation}
\end{comment}

\noindent
Then, from \eqref{s1.1}, \eqref{s2.4},  \eqref{I_J e} and the orthogonality of the ranges of $I_m$ and $J_m$, we have \begin{equation}\label{s2.6}
\left\{
\begin{aligned}
&\partial_t\psi^V=3^{-1}(\nabla\cdot\dot{u}) I_3-b(3\lambda_0+2\mu_0)\psi^V+\sum_{i=1}^n\eta_i^{-1}(3\lambda_i+2\mu_i)\phi_i^V,\\
&\partial_t\phi_i^V=\eta_i^{-1}(3\lambda_0+2\mu_0)\psi^V-\eta_i^{-1}(3\lambda_i+2\mu_i)\phi_i^V,\\
&\psi^V(0)=0=\phi_i^V(0),\quad 1\leq i\leq n
\end{aligned}
\right.
\end{equation}
and
\begin{equation}\label{s2.7}
\left\{
\begin{aligned}
&\partial_t\psi^D=e[\dot{u}]-3^{-1}(\nabla\cdot\dot{u}) I_3-2b\mu_0\psi^D+\sum_{i=1}^n2\mu_i\eta_i^{-1}\phi_i^D,\\
&\partial_t\phi_i^D=2\mu_0\eta_i^{-1}\psi^D-2\mu_i\eta_i^{-1}\phi_i^D,\\
&\psi^D(0)=0=\phi_i^D(0),\quad 1\leq i\leq n.
\end{aligned}
\right.
\end{equation}
Solving \eqref{s2.6} for $\psi^V$, $\phi_i^V$'s, we have 
\begin{equation}\label{s2.8}
\begin{pmatrix}
\psi^V \\
\phi_1^V\\
\cdots  \\
\phi_n^V
\end{pmatrix}
(t)=\int_0^t e^{(t-s)L_0^{\tilde I}}
\begin{pmatrix}
3^{-1}(\nabla\cdot\dot{u})I_3 \\
0\\
\cdots\\
0
\end{pmatrix}
(s)\, ds,
\end{equation}
where
\begin{equation}\label{s2.9}
L_0^{\tilde I}:=
\begin{pmatrix}
                       -b(3\lambda_0+2\mu_0)\tilde I
                       & \eta_1^{-1}(3\lambda_1+2\mu_1)\tilde I & \cdots               &    \eta_n^{-1}(3\lambda_n+2\mu_n) \tilde I   \\

                             \eta_1^{-1}(3\lambda_0+2\mu_0)\tilde I                              & -\eta_1^{-1}(3\lambda_1+2\mu_1)\tilde I & \cdots                  &            0         \\
            \cdots                                   & \cdots                   & \cdots     & \cdots                            \\
                               \eta_{n-1}^{-1}(3\lambda_0+2\mu_0) \tilde I                               & 0              & \cdots          &            0         \\
                              \eta_n^{-1}(3\lambda_0+2\mu_0)\tilde I                              & 0              & \cdots                    &   -\eta_n^{-1}(3\lambda_n+2\mu_n)\tilde I
\end{pmatrix}.
\end{equation}
\noindent
For notational simplicity, we will suppress $\tilde I$ inside the matrix of \eqref{s2.9}, and denote the suppressed one as $L_0$. We remark that the relation between the two is not merely notational, but also a mathematical relation. In fact, consider an $\left(3(n+1)\right)^2$-dimensional Hermitian space $E$ defined by
$$
E:=\{\text{tensor block matrix}\,(\alpha_{ij}\tilde I)_{0\le i,j\le n}:\alpha_{ij}\in{\mathbb R},\,0\le i,\,j\le n\}
$$
with the canonical orthogonal basis consisting of tensor block matrices $\mathcal{E}_{ij}$ for $0\le i,j\le n$ such that each $\mathcal{E}_{ij}$'s $(k,l)$-tensor block component is $\beta_{kl}\tilde I$ with 
$$
\beta_{kl}=\left\{
\begin{array}{ll}
\delta_{ij}\,\,&\text{if}\,\,(k,l)
=(i,j),\\
0\,\,&\text{if otherwise}.
\end{array}
\right.
$$
We note here that the inner product for tensor block matrices is naturally defined as follows. It is the summation of tensor block componentwise inner products, and each of these inner products is the usual componentwise inner product for tensors. 
Then, $L_0^{\tilde I}$ and $L_0$ are algebraically and topologically equivalent.
Similarly, from \eqref{s2.7}, we have
\begin{equation}\label{s2.10}
\begin{pmatrix}
\psi^D \\
\phi^D
\end{pmatrix}
(t)=\int_0^t e^{(t-s)L_1^{\tilde I}}
\begin{pmatrix}
e[\dot{u}]-3^{-1}(\nabla\cdot\dot{u}) I_3 \\
0
\end{pmatrix}
(s)\, ds,
\end{equation}
where
\begin{equation}\label{s2.11}
L_1^{\tilde I}:=
\begin{pmatrix}
                       -2b\mu_0\tilde I
                       & \eta_1^{-1}2\mu_1 \tilde I & \cdots     & \eta_{n-1}^{-1}2\mu_{n-1}  \tilde I          &    \eta_n^{-1}2\mu_n \tilde I   \\

                             \eta_1^{-1}2\mu_0  \tilde I                            & -\eta_1^{-1}2\mu_1 \tilde I& \cdots     &    0               &            0         \\
            \cdots                                   & \cdots                   & \cdots     & \cdots             & \cdots               \\
                               \eta_{n-1}^{-1}2\mu_0    \tilde I                           & 0              & \cdots     & -\eta_{n-1}^{-1}2\mu_{n-1}   \tilde I  &            0         \\
                              \eta_n^{-1}2\mu_0      \tilde I                        & 0              & \cdots     &    0               &    -\eta_n^{-1}2\mu_n \tilde I
\end{pmatrix}.
\end{equation}

Next, we aim to compute $g_{00}(t)$ and analyze its property of the tensor block matrix $e^{tL_1^{I_3}}$ given as
\begin{equation}\label{s2.12}
e^{tL_1^{\tilde I}}=
\begin{pmatrix}
g_{00}(t)\tilde I&g_{01}(t)\tilde I&\cdots&g_{0n}(t)\tilde I\\
\ldots& & &\ldots\\
g_{n0}(t)\tilde I&g_{n1}(t)\tilde I&\cdots&g_{nn}(t)\tilde I
\end{pmatrix}.
\end{equation}
For that consider 
\begin{equation}\label{s2.13}
L_1:=
\begin{pmatrix}
                       -2b\mu_0
                       & \eta_1^{-1}2\mu_1 & \cdots     & \eta_{n-1}^{-1}2\mu_{n-1}            &    \eta_n^{-1}2\mu_n    \\

                             \eta_1^{-1}2\mu_0                              & -\eta_1^{-1}2\mu_1 & \cdots     &    0               &            0         \\
            \cdots                                   & \cdots                   & \cdots     & \cdots             & \cdots               \\
                               \eta_{n-1}^{-1}2\mu_0                               & 0              & \cdots     & -\eta_{n-1}^{-1}2\mu_{n-1}     &            0         \\
                              \eta_n^{-1}2\mu_0                              & 0              & \cdots     &    0               &    -\eta_n^{-1}2\mu_n
\end{pmatrix},
\end{equation}
and compute the $(1,1)$ component of $e^{t L_1}$. Then, it is enough to compute the $(1,1)$ component of the tensor block matrix $e^{t L_1^S}$ and multiply it to $\tilde I$, where $L_1^S$ is given as

\begin{equation}\label{s2.14}
\begin{array}{ll}
L^S_1:=\\
\\
\begin{pmatrix}
                       -b2\mu_0
                       &\eta_1^{-1}\left((2\mu_0)(2\mu_1)\right)^{1/2} & \cdots     & \eta_{n-1}^{-1}\left((2\mu_0)(2\mu_{n-1})\right)^{1/2}            &    \eta_n^{-1}\left((2\mu_0)(2\mu_{n})\right)^{1/2}    \\

                             \eta_1^{-1}\left((2\mu_0)(2\mu_1)\right)^{1/2}                             & -\eta_1^{-1}2\mu_1 & \cdots     &    0               &            0         \\
            \cdots                                   & \cdots                   \cdots     & \cdots             & \cdots               \\
                               \eta_{n-1}^{-1}\left((2\mu_0)(2\mu_{n-1})\right)^{1/2}                                & 0              & \cdots     & -\eta_{n-1}^{-1}2\mu_{n-1}     &            0         \\
                             \eta_n^{-1}\left((2\mu_0)(2\mu_{n})\right)^{1/2}                               & 0              & \cdots     &    0               &   - \eta_n^{-1}2\mu_n
\end{pmatrix}.
\end{array}
\end{equation}


\begin{comment}

Combining \eqref{f8.8} and \eqref{f8.12}, we get that
\begin{equation}\label{f8.13}
\begin{aligned}
\sigma=C_0\psi^S=&\frac{1}{2}\int_0^tC_0e^{-(t-s)b_1}e[\dot{u}](x,s)ds+\frac{1}{2}\int_0^tC_0e^{-(t-s)b_1}e[\dot{u}](x,s)ds\\
=&C_0e[u](x,t)-\frac{b_1}{2}\int_0^tC_0e^{-(t-s)b_1}e[u](x,s)ds-\frac{b_2}{2}\int_0^tC_0e^{-(t-s)b_2}e[u](x,s)ds.
\end{aligned}
\end{equation}
\end{comment}

\noindent
\begin{lemma}\label{s_lem2.1}
$L_1^S$ is negative definite.
\end{lemma}
\begin{proof}
Let $Y=(y_0,y_1,\cdots,y_n)^t$. Then, direct computations of the inner product $(L_1^SY,Y)$ give that
\begin{equation}\label{s2.15}
\begin{aligned}
(L_1^SY,Y)=&-\eta_02\mu_0y^2_0-\sum_{i=1}^n\eta_i^{-1}(\sqrt{2\mu_0}y_0-\sqrt{2\mu_i}y_i)^2\\
\leq &-c|Y|^2
\end{aligned}
\end{equation}
for some constant $c>0$.

\end{proof}

\begin{lemma}\label{s_lem2.2}
$g_{00}(t)$ is the $(1,1)$ element of $e^{tL_1^S}$. 
\end{lemma}
\begin{proof}
Define $A^S$ by
\begin{equation}\label{s2.16}
\begin{aligned}
A^S:=
\begin{pmatrix}
                       -b
                       & \eta_1^{-1} & \cdots     & \eta_{n-1}^{-1}           &    \eta_n^{-1}  \\

                             \eta_1^{-1}                           & -\eta_1^{-1} & \cdots     &    0               &            0         \\
            \cdots                                   & \cdots                   & \cdots     & \cdots             & \cdots               \\
                               \eta_{n-1}^{-1}                                & 0              & \cdots     & -\eta_{n-1}^{-1}     &            0         \\
                             \eta_n^{-1}                               & 0              & \cdots     &    0               &    -\eta_n^{-1}
\end{pmatrix}.
\end{aligned}
\end{equation}
Also, let $D^{\mu}={\rm diag}(\sqrt{2\mu_0},\sqrt{2\mu_1},\cdots,\sqrt{2\mu_n})$. Then, it is easy to see that 
\begin{equation}\label{s2.17}
\begin{aligned}
D^{\mu}A^SD^{\mu}=L_1^S.
\end{aligned}
\end{equation}
Now, observe that $g_{00}(t)$ is the $(1,1)$ element of $(2\mu_0)^{-1}D^{\mu}D^{\mu}e^{tL_1}$. Further, since $L_1=A^SD^{\mu}D^{\mu}$, we have 
\begin{equation}\label{s2.18}
\begin{aligned}
D^{\mu}D^{\mu}e^{tL_1}=D^{\mu}e^{tD^{\mu}A^SD^{\mu}}D^{\mu}=D^{\mu}e^{tL_1^S}D^{\mu}.
\end{aligned}
\end{equation}
Then, from \eqref{s2.18}, we have that $g_{00}(t)$ is the $(1,1)$ element of $e^{tL_1^S}$.

\end{proof}

\noindent
Based on Lemma \ref{s_lem2.1}, let $-\tau_j<0$ for $0\leq j\leq n$ be the eigenvalues of $L_1^S$ and the associated unit eigenvectors $v^j$ for $0\leq j\leq n$ such that
$L_1^Sv^j=-\tau_j v^j$ for $0\leq j\leq n$.
Also, let
$D^{\tau}={\rm diag}(-\tau_0,-\tau_1,\cdots,-\tau_n)$ and 
$V$ be a matrix with $n+1$-dimensional unit column vectors $v^j,\,0\le j\le n$ each of which has the i-th component $v_i^j$. Then we have
\begin{equation}\label{s2.19}
\begin{aligned}
L_1^S=VD^{\tau}V^t,
\end{aligned}
\end{equation}
which yields the following theorem.
\begin{Th}\label{s_thm2.3}
We have the following expression for $g_{00}(t)$:
\begin{equation}\label{s2.20}
\begin{aligned}
g_{00}(t)=\sum_{j=0}^n(v_0^j)^2e^{-t\tau_j}.
\end{aligned}
\end{equation}
\end{Th}
\begin{proof}
By using \eqref{s2.19}, we have that
\begin{equation}\label{s2.21}
\begin{aligned}
e^{tL_1^S}=e^{tVD^{\tau}V^t}=Ve^{tD^{\tau}}V^t.
\end{aligned}
\end{equation}
Combining Lemma \ref{s_lem2.2} and \eqref{s2.21}, we obtain \eqref{s2.20}.

\end{proof}


Next, in the same way as we did for $e^{t L_1^{\tilde I}}$, we compute the $(1,1)$ tensor block component of the tensor block matrix $e^{t L_0^{\tilde I}}$. For that let
\begin{equation}\label{s2.22}
L_0^S:=
\begin{pmatrix}
                       -b(3\lambda_0+2\mu_0)
                       & \eta_1^{-1}h_1  & \cdots    &        \eta_{n-1}^{-1}h_{n-1}     &    \eta_n^{-1}h_n   \\

                              \eta_1^{-1}h_1                             & -\eta_1^{-1}(3\lambda_1+2\mu_1) & \cdots                  &          0  &   0         \\
            \cdots                                   & \cdots                   & \cdots     & \cdots           & \cdots                 \\
                               \eta_{n-1}^{-1}h_{n-1}                                & 0              & \cdots          &      -\eta_{n-1}^{-1}(3\lambda_{n-1}+2\mu_{n-1}) &      0         \\
                              \eta_n^{-1}h_n                              & 0              & \cdots        & 0            &   -\eta_n^{-1}(3\lambda_n+2\mu_n)
\end{pmatrix},
\end{equation}
where $h_i=\sqrt{(3\lambda_0+2\mu_0)(3\lambda_i+2\mu_i)} $ for $1\leq i \leq n$.

\medskip
\noindent
Then, we can prove the following lemma in the same as we did for $L_1^S$.

\begin{lemma}\label{s_lem2.4}
$L_0^S$ is negative definite.
\end{lemma}

\medskip
\noindent
Further, by representing
$e^{tL_0(x)}$ as
\begin{equation}\label{s2.23}
e^{tL_0}=
\begin{pmatrix}
g^0_{00}(t)&g^0_{01}(t)&\cdots&g^0_{0n}(t)\\
\ldots& & &\ldots\\
g^0_{n0}(t)&g^0_{n1}(t)&\cdots&g^0_{nn}(t)
\end{pmatrix},
\end{equation}
we can prove the following lemma in the same way as we proved Lemma \ref{s_lem2.2}.
\begin{lemma}\label{s_lem2.5}
$g^0_{00}(t)$ is the $(1,1)$ element of $e^{tL_0^S}$. 
\end{lemma}

\medskip
\noindent
Based on Lemma \ref{s_lem2.4}, let $-\kappa_j<0$ for $0\leq j\leq n$ be the eigenvalues of $L_0^S$ with the associated unit eigenvectors $q^j$ for $0\leq j\leq n$ such that
$L_0^Sq^j=-\kappa_j q^j$ for $0\leq j\leq n$. Also, let
$D^{\kappa}={\rm diag}(-\kappa_0,-\kappa_1,\cdots,-\kappa_n)$ and $Q$ be a matrix with $n+1$-dimensional unit column vectors $q^j,\,0\le j\le n$ each of which has the i-th component $q_i^j$. Then we have
\begin{equation}\label{s2.24}
\begin{aligned}
L_0^S=QD^{\kappa}Q^{\mathfrak{t}},
\end{aligned}
\end{equation}
where $\mathfrak{t}$ denotes the transpose.
\medskip\noindent
Then, we can prove that this yields the following theorem in the same way as we proved Theorem \ref{s_thm2.3}

\begin{Th}\label{s_thm2.6}
We have the following expression for $g_{00}^0(t)$
\begin{equation}\label{s2.25}
\begin{aligned}
g^0_{00}(t)=\sum_{j=0}^n(q_0^j)^2e^{-t\kappa_j}.
\end{aligned}
\end{equation}
\end{Th}

Now, combining \eqref{s2.3}, \eqref{s2.20} and \eqref{s2.25}, we have 
\begin{equation}\label{s2.26}
\begin{aligned}
\psi=\psi^V+\psi^D=&\int_0^t g^0_{00}(t-s)3^{-1}(\nabla\cdot\dot{u})(s)\,ds\,I_3\\
&+\int_0^t g_{00}(t-s)\big(e[\dot{u}]-3^{-1}(\nabla\cdot\dot{u})I_3\big) (s)\,ds\\
=&\int_0^t \sum_{j=0}^n(q_0^j)^2e^{-(t-s)\kappa_j}3^{-1}(\nabla\cdot\dot{u})(s)\,ds\,I_3\\
&+\int_0^t \sum_{j=0}^n(v_0^j)^2e^{-(t-s)\tau_j}\big(e[\dot{u}]-3^{-1}(\nabla\cdot\dot{u})I_3\big) (s)\,ds.
\end{aligned}
\end{equation}
Further, from \eqref{s2.26}, we have the following explicit form of the stress-strain relation, which yields the explicit form of the relaxation tensor for the inhomogeneous isotropic EBM:
\begin{equation}\label{s2.27}
\begin{aligned}
\sigma=C_0\psi
=&(\lambda_0+\frac{2}{3}\mu_0)\int_0^t \sum_{j=0}^n(q_0^j)^2e^{-(t-s)\kappa_j}(\nabla\cdot\dot{u})(s)\,ds\,I_3\\
&+2\mu_0\int_0^t \sum_{j=0}^n(v_0^j)^2e^{-(t-s)\tau_j}\big(e[\dot{u}]-3^{-1}(\nabla\cdot\dot{u})I_3\big) (s)\,ds.
\end{aligned}
\end{equation}
By rewriting this using $e[u](0)=0$, and $\sum_{j=0}^n (q_0^j)^2=\sum_{j=0}^n (v_0^j)^2=1$, we have 
\begin{equation}\label{s2.28}
\begin{aligned}
\sigma
=&\big((\lambda_0+\frac{2}{3}\mu_0)(\nabla\cdot u)I_3\big)(t)+ 2\mu_0\big(e[u]-3^{-1}(\nabla\cdot u)I_3\big) (t)\\
&-(\lambda_0+\frac{2}{3}\mu_0)\int_0^t \sum_{j=0}^n\kappa_j(q_0^j)^2e^{-(t-s)\kappa_j}(\nabla\cdot u)(s)\,ds\,I_3\\
&-2\mu_0\int_0^t \sum_{j=0}^n\tau_j(v_0^j)^2e^{-(t-s)\tau_j}\big(e[u]-3^{-1}(\nabla\cdot u)I_3\big) (s)\,ds.
\end{aligned}
\end{equation}


\begin{remark}
\begin{itemize}

\noindent
\item[\rm{(i)}]
Up to scalar multiplication,  $(\nabla\cdot\dot u)I_3$ and $e[\dot u]-3^{-1}(\nabla\cdot\dot u)I_3$ in \eqref{s2.27} correspond to the volumetric part and deviatoric part, respectively. They are orthogonal. However, their tractions, the inner product of their contractions with the unit outer normal vector $\nu$ of $\partial\Omega$, are given as
$$
(\nabla\cdot\dot u)(e[\dot u],\nu\otimes\nu)-3^{-1}(\nabla\cdot\dot u)^2
$$
does not necessarily vanish, where $(\,\,,\,\,)$ denotes the inner product for the second order tensors. Hence, concerning the traction-free boundary condition, it is very hard to split them.
\item[\rm{(ii)}] In \eqref{s2.28}, $(\nabla\cdot u)I_3$ and $e[u]-3^{-1}(\nabla\cdot u)I_3$ appear jointly and separately outside and inside the integral, respectively. This is due to the new properties of the relaxation tensor that arise from the two aforementioned parts, and we present a new idea in the two forthcoming sections to generate the C-ev's.

\end{itemize}
\end{remark}

 

\begin{comment}
\section{Burger models and special models}\label{burger}
\setcounter{equation}{0}

For the case $n=1$, $D^{\tau}={\rm diag}(-\tau_0,-\tau_1)$, where $$\tau_0=\mu_0b+\mu_1\eta_1^{-1}+\sqrt{(\mu_0b-\mu_1\eta_1^{-1})^2+4\mu_0\mu_1\eta^{-2}}$$
and
$$\tau_1=\mu_0b+\mu_1\eta_1^{-1}-\sqrt{(\mu_0b-\mu_1\eta_1^{-1})^2+4\mu_0\mu_1\eta^{-2}}.$$
Also, we have
\begin{equation}\label{I6.1}
\begin{aligned}
g_{00}(x,t)=\frac{4\mu_0\mu_1\eta_1^{-2}}{4\mu_0\mu_1\eta_1^{-2}+(-\tau_0+2\mu_0b)^2}e^{t\tau_0}+\frac{4\mu_0\mu_1\eta_1^{-2}}{4\mu_0\mu_1\eta_1^{-2}+(-\tau_1+2\mu_0b)^2}e^{t\tau_1}.
\end{aligned}
\end{equation}
$D^{\kappa}={\rm diag}(-\kappa_0,-\kappa_1)$, where $$\kappa_0=\frac{b(d\lambda_0+2\mu_0)+\eta_1^{-1}(d\lambda_1+2\mu_1)}{2}+\frac{1}{2}\sqrt{[b(d\lambda_0+2\mu_0)-\eta_1^{-1}(d\lambda_1+2\mu_1)]^2+4h_1\eta^{-2}}$$
and
$$\kappa_1=\frac{b(d\lambda_0+2\mu_0)+\eta_1^{-1}(d\lambda_1+2\mu_1)}{2}-\frac{1}{2}\sqrt{[b(d\lambda_0+2\mu_0)-\eta_1^{-1}(d\lambda_1+2\mu_1)]^2+4h_1\eta^{-2}}.$$
Also, we have
\begin{equation}\label{I6.2}
\begin{aligned}
g^0_{00}(x,t)=\frac{h^2_1\eta_1^{-2}}{h^2_1\eta_1^{-2}+(-\kappa_0+bd\lambda_0+2b\mu_0)^2}e^{t\kappa_0}+\frac{h^1_2\eta_1^{-2}}{h^2_1\eta_1^{-2}+(-\kappa_1+bd\lambda_0+2b\mu_0)^2}e^{t\kappa_1}.
\end{aligned}
\end{equation}
From \eqref{I5.14}, we get that
\begin{equation}\label{I6.3}
\begin{aligned}
\sigma
=&(\lambda_0+2\mu_0d^{-1})\int_0^t \frac{h^2_1\eta_1^{-2}}{h^2_1\eta_1^{-2}+(-\kappa_0+bd\lambda_0+2b\mu_0)^2}e^{-(t-s)\kappa_0}(\nabla\cdot\dot{u})(x,s)\,ds\,I_d\\
&+(\lambda_0+2\mu_0d^{-1})\int_0^t \frac{h^2_1\eta_1^{-2}}{h^2_1\eta_1^{-2}+(-\kappa_1+bd\lambda_0+2b\mu_0)^2}e^{-(t-s)\kappa_1}(\nabla\cdot\dot{u})(x,s)\,ds\,I_d\\
&-2\mu_0d^{-1}\int_0^t \frac{4\mu_0\mu_1\eta_1^{-2}}{4\mu_0\mu_1\eta_1^{-2}+(-\tau_0+2\mu_0b)^2}e^{-(t-s)\tau_0}(\nabla\cdot\dot{u}) (x,s)\,ds\,I_d\\
&-2\mu_0d^{-1}\int_0^t \frac{4\mu_0\mu_1\eta_1^{-2}}{4\mu_0\mu_1\eta_1^{-2}+(-\tau_1+2\mu_0b)^2}e^{-(t-s)\tau_1}(\nabla\cdot\dot{u}) (x,s)\,ds\,I_d\\
&+ 2\mu_0\int_0^t \frac{4\mu_0\mu_1\eta_1^{-2}}{4\mu_0\mu_1\eta_1^{-2}+(-\tau_0+2\mu_0b)^2}e^{-(t-s)\tau_0}e[\dot{u}] (x,s)\,ds\\
&+ 2\mu_0\int_0^t \frac{4\mu_0\mu_1\eta_1^{-2}}{4\mu_0\mu_1\eta_1^{-2}+(-\tau_1+2\mu_0b)^2}e^{-(t-s)\tau_1}e[\dot{u}] (x,s)\,ds.
\end{aligned}
\end{equation}

By Theorem \ref{B_thm7.4}
\begin{Th}\label{I_thm6.1}
Assume that $b\mu_0=\eta_j^{-1}\mu_j$ for $1\leq j\leq n$. Then
\begin{equation}\label{I6.4}
\begin{aligned}
g_{00}(x,t)=&\frac{1}{2}\{e^{-t(2b\mu_0+2b\mu_0\sqrt{b^{-1}\sum_{i=1}^n\eta_i^{-1}}\,)}+e^{-t(2b\mu_0-2b\mu_0\sqrt{b^{-1}\sum_{i=1}^n\eta_i^{-1}}\,)}\}\\
=&\frac{1}{2}e^{-t\tau_0}+\frac{1}{2}e^{-t\tau_1},
\end{aligned}
\end{equation}
where $b=\sum_{i=0}^n\eta_i^{-1}$, $\tau_0=2b\mu_0(1+\sqrt{b^{-1}(\sum_{i=1}^n\eta_i^{-1})})$ and 
$\tau_1=2b\mu_0(1-\sqrt{b^{-1}(\sum_{i=1}^n\eta_i^{-1})})$.
\end{Th}

\begin{Th}\label{I_thm6.2}
Assume that $b\mu_0=\eta_j^{-1}\mu_j$ and $b\lambda_0=\eta_j^{-1}\lambda_j$ for $1\leq j\leq n$. Then we have \eqref{I6.4} and
\begin{equation}\label{I6.5}
\begin{aligned}
g^0_{00}(x,t)=\frac{1}{2}e^{t\kappa_0}+\frac{1}{2}e^{t\kappa_1},
\end{aligned}
\end{equation}
where $\kappa_0=b(d\lambda_0+2\mu_0)(1+\sqrt{b^{-1}(\sum_{i=1}^n\eta_i^{-1})})$ and 
$\kappa_1=b(d\lambda_0+2\mu_0)(1-\sqrt{b^{-1}(\sum_{i=1}^n\eta_i^{-1})})$. Moreover, we have
\begin{equation}\label{I6.6}
\begin{aligned}
\sigma
=&\int_0^t[\frac{(d\lambda_0+2\mu_0)}{2d} (e^{-(t-s)\kappa_0}+e^{-(t-s)\kappa_1})-\mu_0d^{-1}(e^{-(t-s)\tau_0}+e^{-(t-s)\tau_1})](\nabla\cdot\dot{u})(x,s)\,ds\,I_d\\
&+ \mu_0\int_0^t (e^{-(t-s)\tau_0}+e^{-(t-s)\tau_1})e[\dot{u}] (x,s)\,ds.
\end{aligned}
\end{equation}
\end{Th}




\section{Eigenvalues for homogeneous isotropic elasticity}\label{cluster eigenvalues}
\setcounter{equation}{0}
From \eqref{I5.15}, we have for homogeneous isotropic elasticity that

\begin{equation}\label{I7.1}
\begin{aligned}
\nabla\cdot\sigma
=&(\lambda_0+\mu_0)\nabla p(x,t)+ \mu_0\Delta u(x,t)\\
&+(\lambda_0+2\mu_0d^{-1})\int_0^t \sum_{j=0}^n\kappa_j(q_0^j)^2e^{-(t-s)\kappa_j}\nabla p(x,s)\,ds\\
&+\mu_0(1-2d^{-1})\int_0^t \sum_{j=0}^n\tau_j(v_0^j)^2e^{-(t-s)\tau_j}\nabla p(x,s)\,ds\\
&+ \mu_0\int_0^t \sum_{j=0}^n\tau_j(v_0^j)^2e^{-(t-s)\tau_j}\Delta u (x,s)\,ds,
\end{aligned}
\end{equation}
where $p=\nabla \cdot u$.
Taking divergence on
$$\partial_{tt}u=\nabla\cdot\sigma,$$
we get that
\begin{equation}\label{I7.2}
\begin{aligned}
\partial_{tt} p(x,t)
=&(\lambda_0+2\mu_0)\Delta p(x,t)\\
&+(\lambda_0+2\mu_0d^{-1})\int_0^t \sum_{j=0}^n\kappa_j(q_0^j)^2e^{-(t-s)\kappa_j}\Delta p(x,s)\,ds\\
&+\mu_0(2-2d^{-1})\int_0^t \sum_{j=0}^n\tau_j(v_0^j)^2e^{-(t-s)\tau_j}\Delta p(x,s)\,ds.
\end{aligned}
\end{equation}
\begin{lemma}
Let $u$ satisfy
\label{I_lem7.1}
\begin{equation}\label{I7.3}
\begin{aligned}
\nabla \cdot (C_0e[u])=k_u\, u\quad &\mbox{\rm in}\quad \Omega
\end{aligned}
\end{equation}  
for some constant $k_u$. 
Then we have
\begin{equation}\label{I7.4}
\begin{aligned}\Delta p(x)=k_p\, p(x),
\end{aligned}
\end{equation}
where $k_p=(\lambda_0+2\mu_0)^{-1}k_u$ and $p=\nabla \cdot u$.
\end{lemma}

\begin{proof}
Taking divergence on the first equation of \eqref{I7.3}, we have that
\begin{equation}\label{I7.5}
\begin{aligned}
\Delta p(x)=(\lambda_0+2\mu_0)^{-1}k_u\,  p(x)=k_p\, p(x),
\end{aligned}
\end{equation}
where $k_p=(\lambda_0+2\mu_0)^{-1}k_u$.
\end{proof}

On the other hand, we have the following Lemma.
\begin{lemma}
Let  $p$ satisfy
\begin{equation}\label{I7.6}
\begin{aligned}
\Delta p(x)=k_p\, p(x),
\end{aligned}
\end{equation}
where $k_p$ is a constant. 
Then we have
\begin{equation}\label{I7.7}
\begin{aligned}
\nabla \cdot (C_0e[u])=k_u\, u,
\end{aligned}
\end{equation}
where $k_u=(\lambda_0+2\mu_0)k_p$ and $u=\nabla p$.
\end{lemma}

\begin{proof}
The direct computation gives that
\begin{equation}\label{I7.8}
\begin{aligned}
\nabla \cdot (C_0e[u])&=(\lambda_0+\mu_0)\nabla (\nabla\cdot u)+\mu_0\Delta u \\
&=(\lambda_0+\mu_0)k_p\nabla p+\mu_0k_p \nabla p\\
&=(\lambda_0+2\mu_0)k_p\, u\\
&=k_u\, u.
\end{aligned}
\end{equation}
\end{proof}


For simplicity, let $\xi p:=\Delta p $ and $I(r):=\int_0^t e^{-r(t-\tau)} \Delta p(\tau,x) \, d\tau $. The previous method for \eqref{I7.2} should work.


\section{Cluster eigenvalues for homogeneous isotropic elasticity2}\label{cluster eigenvalues2}
\setcounter{equation}{0}

Recall \eqref{I5.15}, we have that
\begin{equation*}
\begin{aligned}
\sigma
=&\lambda_0(\nabla\cdot u)(x,t)\,I_d+ 2\mu_0e[u] (x,t)\\
&+(\lambda_0+2\mu_0d^{-1})\int_0^t \sum_{j=0}^n\kappa_j(q_0^j)^2e^{-(t-s)\kappa_j}(\nabla\cdot u)(x,s)\,ds\,I_d\\
&-2\mu_0d^{-1}\int_0^t \sum_{j=0}^n\tau_j(v_0^j)^2e^{-(t-s)\tau_j}(\nabla\cdot u) (x,s)\,ds\,I_d\\
&+ 2\mu_0\int_0^t \sum_{j=0}^n\tau_j(v_0^j)^2e^{-(t-s)\tau_j}e[u] (x,s)\,ds.
\end{aligned}
\end{equation*}
Let us define two operators.
\begin{equation}\label{I8.1}
\begin{aligned}
Q_A(x)u=\nabla\cdot \big(d^{-1}(\nabla\cdot u)I_d\big)
\end{aligned}
\end{equation}
and
\begin{equation}\label{I8.2}
\begin{aligned}
Q_B(x)u=\nabla\cdot \big(e[u]-d^{-1}(\nabla\cdot u)I_d\big).
\end{aligned}
\end{equation}
\begin{lemma}\label{I_lem8.1}
Let $u$ satisfy $Q_A(x)u=k_a\, u$ and $Q_B(x)u=(1-d^{-1})dk_a\, u$ for some constant $k_a$. Then we have
$$\Delta p(x)=k_b\, p(x),$$
where $k_b=dk_a$ and $p=\nabla \cdot u$.
\end{lemma}
\begin{proof}
The direct computations give that
\begin{equation}\label{I8.3}
\begin{aligned}
 \nabla \cdot (Q_A(x)u)=d^{-1}\Delta p(x)=k_a\, p
\end{aligned}
\end{equation}
and
\begin{equation}\label{I8.4}
\begin{aligned}
 \nabla \cdot (Q_B(x)u)=(1-d^{-1})dk_a\, p.
\end{aligned}
\end{equation}
\end{proof}

\begin{lemma}\label{I_lem8.2}
Let $\Delta p(x)=k_b\, p(x)$ and $u=\nabla p$. Then $Q_A(x)u=d^{-1}k_b\, u$ and $Q_B(x)u=(1-d^{-1})k_b\, u$.
\end{lemma}

\end{comment}

\begin{comment}
    
\section{Eigenvalues of Lam\'e operator {\color{magenta} in charge of Youjun}}
In this section, we shall consider the eigenvalues to the Lam\'e operator $\nabla \cdot C_0e$, where $e[u]$ represent the strain tensor. Denote by $j_n$ and $h_n$ the spherical bessel and spherical Hankel function of first order and degree $n$, respectively. It holds that
\[
h_n(t)=j_n(t)+iy_n(t).
\]
Recall that, in \cite{kong,spectraltan}, three types of vectorial polynomials are introduced:
\begin{equation}
\mathcal{T}^m_n(x)=\nabla (r^{n} Y^m_{n}(\hat{x}))\times x, \quad n\geqslant 1, \quad  -n\leqslant m\leqslant n,\nonumber
\end{equation}
\[
\mathcal{M}^m_n(x)=\nabla(r^{n} Y^m_{n}(\hat{x})), \quad n\geqslant 1,\quad  -n\leqslant m\leqslant n,
\]
and
\[
\mathcal{N}^m_n(x)=a^m_n r^{n-1} Y^m_{n-1}(\hat{x})x + (1-\frac{a_n^m}{2n-1}-r^2) \nabla(r^{n-1} Y^m_{n-1}(\hat{x})),
\]
where
\begin{equation}\label{eq:43}
a^m_n = \frac{2(n-1)\lambda_0 + 2(3n-2)\mu_0}{(n+2)\lambda_0 + (n+4)\mu_0}, \quad n\geqslant 1, \quad  -(n-1)\leqslant m\leqslant n-1,
\end{equation}
and $\hat{x}=x/|x|$,
which are solutions to $\nabla \cdot(C_0e[u]) =0$ in $\mathbb{R}^3$. The traces of $\mathcal{ T}^m_n,\mathcal{M}^m_n,\mathcal{N}^m_n$ on the unit sphere $S$ denoted by $\mathbf{T}^m_n,\mathbf{M}^m_n,\mathbf{N}^m_n$, have the form
\begin{equation}
\mathbf{T}_n^m(x)= \nabla_{S}Y_{n}^m(\hat{x})\times\nu, \nonumber%\quad n\geqslant 1, \quad  -n\leqslant m\leqslant n,\nonumber
\end{equation}
\begin{equation}
\mathbf{M}_n^m(x)= \nabla_{S}Y_{n}^m(\hat{x}) +n Y_{n}^m(\hat{x}) \nu,\nonumber% \quad n\geqslant 1,\quad  -n\leqslant m\leqslant n,\nonumber
\end{equation}
\begin{equation}
\mathbf{N}_n^m(x)= \frac{a^m_n}{2n-1}(-\nabla_{S}Y_{n-1}^m(\hat{x}) +n Y_{n-1}^m(\hat{x})\nu), \nonumber%\quad n\geqslant 1,\quad  -n\leqslant m\leqslant n,\nonumber
\end{equation}%By straight forward computations  one can also verify that
which form an orthogonal basis of $L^2(S)^3$. 
For later usage, we also define
\[
\tilde{\mathbf{N}}_n^m(x)=\frac{2n-1}{a_n^m}\mathbf{N}_n^m(x).
\]

%\begin{comment}

\subsection{Dirichlet boundary condition}
Let us first assume $\Gamma_N=\emptyset$, then the boundary condition appeared in (\ref{B2.4}) is only Dirichlet boundary condition. We consider the eigenvalues to the following system
\begin{equation}\label{I9.1}
\left\{\begin{aligned}
\nabla \cdot (C_0e[u])=-\omega^2\, u\quad &\mbox{\rm in}\quad \Omega \\
u=0, \quad &\mbox{\rm on}\quad \partial \Omega.
\end{aligned}
\right.
\end{equation}  
Suppose that $\Omega$ is a unit ball. 

Introducing the fundamental solution $\mathbf{\Gamma}^{\omega}$ to $\nabla \cdot C_0e +\omega^2$ in $\mathbb{R}^3$ is expressed as
\begin{equation}\label{eq:33}
\mathbf{\Gamma}^{\omega}(x)=-\frac{e^{ik_s|x|}}{4\pi \mu_0|x|}\mathbf{I}+\frac{1}{4\pi\omega^{2}\rho}\nabla\nabla(\frac{e^{ik_{p}|x|}-e^{ik_{s}|x|}}{|x|}),%\nonumber
\end{equation}
where $k_{s}=\frac{\omega}{c_{s}},k_{p}=\frac{\omega}{c_{p}}$ with
\begin{equation}\label{eq:17}
c_{s}=\sqrt{\mu_0},\quad c_{p}=\sqrt{\lambda_0+2\mu_0}.
\end{equation}
Specially, when $\omega=0$, the fundamental solution to $\nabla \cdot C_0e$ is given by
\begin{equation}\label{eq:35}
\mathbf{\Gamma}^{0}(x)=-\frac{\hat{\gamma}_{1}}{4\pi}\frac{1}{|x|}\mathbf{I}-\frac{\hat{\gamma}_{2}}{4\pi}\frac{xx^{T}}{|x|^{3}},%\nonumber
\end{equation}
where $\hat{\gamma}_{1}=\frac{1}{2}(\frac{1}{\mu_0}+\frac{1}{2\mu)0+\lambda_0}),\hat{\gamma}_{2}=\frac{1}{2}(\frac{1}{\mu_0}-\frac{1}{2\mu_0+\lambda_0})$.
Based on the above notation, for any $\varphi\in H^{-1/2}(\partial \Omega)^{3}$, the single layer potential can be written by
\begin{equation}
\mathbf{S}^\omega_\Omega[ \varphi](x):=\int_{\partial \Omega}\mathbf{\Gamma}^{\omega}(x-y)\varphi(y)d\sigma_{y},\quad x\in \mathbb{R}^3.\nonumber
\end{equation}

Recall the following elementary result (cf.\cite{DLL}):
\[
\begin{split}
  \mathbf{S}^\omega_\Omega[ \mathbf{M}_{n}^m] =&  -i \left( \frac{ (n+1)k_s j_{n-1}(k_s) h_{n-1}(k_s) }{\mu_0(2n+1)}  +  \frac{ nk_p j_{n-1}(k_p) h_{n-1}(k_p)}{(\lambda_0+2\mu_0)(2n+1)}   \right)\mathbf{M}_{n}^m \\
       &- i n    \left( \frac{k_s j_{n-1}(k_s) h_{n+1}(k_s) }{\mu(2n+1)}  -   \frac{k_p j_{n-1}(k_p)h_{n+1}(k_p ) }{(\lambda_0+2\mu_0)(2n+1)}  \right) \tilde{\mathbf{N}}_{n+1}^m,
 \end{split}
\]

\[
\begin{split}
  \mathbf{S}^\omega_\Omega[ \tilde{\mathbf{N}}_{n+1}^m] =&  - i(n+1)   \left( \frac{ k_s j_{n+1}(k_s) h_{n-1}(k_s)}{\mu_0(2n+1)}  -  \frac{ k_p j_{n+1}(k_p) h_{n-1}(k_p) }{(\lambda_0+2\mu_0)(2n+1)}   \right) \mathbf{M}_{n}^m \\
       & -i \left( \frac{ n k_s j_{n+1}(k_s) h_{n+1}(k_s) }{\mu_0(2n+1)}  +  \frac{ (n+1) k_p j_{n+1}(k_p) h_{n+1}(k_p )}{(\lambda_0+2\mu_)(2n+1)}  \right) \tilde{\mathbf{N}}_{n+1}^m,
 \end{split}
\]
and 
\[
 \mathbf{S}^\omega_\Omega[\mathbf{T}_n^m] = -\frac{i k_s j_{n}(k_s) h_n(k_s)}{\mu}\mathbf{T}_n^m, 
\]
hold on $\partial \Omega$.

Let us first focus on the real eigenvalues, that is, we assume $\omega\in\mathbb{R}$. There are two different modes to be considered. First, consider the linear combination of  $\mathbf{M}_{n}^m$ and $\mathbf{N}_{n+1}^m$. Define
\[
f_{1n}=\frac{ (n+1)k_s j_{n-1}(k_s) h_{n-1}(k_s) }{\mu_0(2n+1)}  +  \frac{ nk_p j_{n-1}(k_p) h_{n-1}(k_p)}{(\lambda_0+2\mu_0)(2n+1)},
\]
\[
g_{1n}=\frac{k_s j_{n-1}(k_s) h_{n+1}(k_s) }{\mu_0(2n+1)}  -   \frac{k_p j_{n-1}(k_p)h_{n+1}(k_p ) }{(\lambda_0+2\mu_0)(2n+1)},
\]
and
\[
f_{2n}=\frac{ k_s j_{n+1}(k_s) h_{n-1}(k_s)}{\mu_0(2n+1)}  -  \frac{ k_p j_{n+1}(k_p) h_{n-1}(k_p) }{(\lambda_0+2\mu_0)(2n+1)},
\]
\[
g_{2n}=\frac{ n k_s j_{n+1}(k_s) h_{n+1}(k_s) }{\mu_0(2n+1)}  +  \frac{ (n+1) k_p j_{n+1}(k_p) h_{n+1}(k_p )}{(\lambda_0+2\mu_)(2n+1)}.
\]
Then there should satisfy
\[
\left|
\begin{array}{cccc}
   f_{1n}^r &-f_{1n}^i&g_{1n}^r  & -g_{1n}^i \\
   f_{1n}^r &f_{1n}^i&g_{1n}^r  & g_{1n}^i \\
   f_{2n}^r &-f_{2n}^i&g_{2n}^r  & -g_{2n}^i \\
   f_{2n}^r &f_{2n}^i&g_{2n}^r  & g_{2n}^i
\end{array}
\right|=0
\]
where $f_{jn}^r$, $g_{jn}^r$, $j=1,2$ are real parts of $f_{jn}$ and $g_{jn}$, respectively. $f_{jn}^i$, $g_{jn}^i$, $j=1,2$ are imaginary parts of $f_{jn}$ and $g_{jn}$, respectively.


Second, consider $\mathbf{T}_n^m$, then it only requires
\[
j_n(k_s)h_n(k_s)=0,
\]
which means that $k_s$ is a root of $j_n$. And the eigenvalues are $-k_s^2\mu_0^2$.

\subsection{Traction free boundary condition}



Introducing the fundamental solution $\mathbf{\Gamma}^{\omega}$ to $\nabla \cdot C_0e +\omega^2$ in $\mathbb{R}^3$ is expressed as
\begin{equation}\label{eq:33}
\mathbf{\Gamma}^{\omega}(x)=-\frac{e^{ik_s|x|}}{4\pi \mu_0|x|}\mathbf{I}+\frac{1}{4\pi\omega^{2}\rho}\nabla\nabla(\frac{e^{ik_{p}|x|}-e^{ik_{s}|x|}}{|x|}),%\nonumber
\end{equation}
where $k_{s}=\frac{\omega}{c_{s}},k_{p}=\frac{\omega}{c_{p}}$ with
\begin{equation}\label{eq:17}
c_{s}=\sqrt{\mu_0},\quad c_{p}=\sqrt{\lambda_0+2\mu_0}.
\end{equation}
Specially, when $\omega=0$, the fundamental solution to $\nabla \cdot C_0e$ is given by
\begin{equation}\label{eq:35}
\mathbf{\Gamma}^{0}(x)=-\frac{\hat{\gamma}_{1}}{4\pi}\frac{1}{|x|}\mathbf{I}-\frac{\hat{\gamma}_{2}}{4\pi}\frac{xx^{T}}{|x|^{3}},%\nonumber
\end{equation}
where $\hat{\gamma}_{1}=\frac{1}{2}(\frac{1}{\mu_0}+\frac{1}{2\mu)0+\lambda_0}),\hat{\gamma}_{2}=\frac{1}{2}(\frac{1}{\mu_0}-\frac{1}{2\mu_0+\lambda_0})$.
Based on the above notation, for any $\varphi\in H^{-1/2}(\partial \Omega)^{3}$, the single layer potential can be written by
\begin{equation}
\mathbf{S}^\omega_\Omega[ \varphi](x):=\int_{\partial \Omega}\mathbf{\Gamma}^{\omega}(x-y)\varphi(y)d\sigma_{y},\quad x\in \mathbb{R}^3.\nonumber
\end{equation}

Recall the following elementary result (cf.\cite{DLL}):
\[
\begin{split}
  \mathbf{S}^\omega_\Omega[ \mathbf{M}_{n}^m] =&  -i \left( \frac{ (n+1)k_s j_{n-1}(k_s) h_{n-1}(k_s) }{\mu_0(2n+1)}  +  \frac{ nk_p j_{n-1}(k_p) h_{n-1}(k_p)}{(\lambda_0+2\mu_0)(2n+1)}   \right)\mathbf{M}_{n}^m \\
       &- i n    \left( \frac{k_s j_{n-1}(k_s) h_{n+1}(k_s) }{\mu(2n+1)}  -   \frac{k_p j_{n-1}(k_p)h_{n+1}(k_p ) }{(\lambda_0+2\mu_0)(2n+1)}  \right) \tilde{\mathbf{N}}_{n+1}^m,
 \end{split}
\]

\[
\begin{split}
  \mathbf{S}^\omega_\Omega[ \tilde{\mathbf{N}}_{n+1}^m] =&  - i(n+1)   \left( \frac{ k_s j_{n+1}(k_s) h_{n-1}(k_s)}{\mu_0(2n+1)}  -  \frac{ k_p j_{n+1}(k_p) h_{n-1}(k_p) }{(\lambda_0+2\mu_0)(2n+1)}   \right) \mathbf{M}_{n}^m \\
       & -i \left( \frac{ n k_s j_{n+1}(k_s) h_{n+1}(k_s) }{\mu_0(2n+1)}  +  \frac{ (n+1) k_p j_{n+1}(k_p) h_{n+1}(k_p )}{(\lambda_0+2\mu_)(2n+1)}  \right) \tilde{\mathbf{N}}_{n+1}^m,
 \end{split}
\]
and 
\[
 \mathbf{S}^\omega_\Omega[\mathbf{T}_n^m] = -\frac{i k_s j_{n}(k_s) h_n(k_s)}{\mu}\mathbf{T}_n^m, 
\]
hold on $\partial \Omega$.
Assume $\Gamma_D=\emptyset$, then the boundary condition appeared in (\ref{s2.1}) is only Traction free boundary condition. We consider the eigenvalues to the following system
\begin{equation}\label{I9.10}
\left\{\begin{aligned}
\nabla \cdot (C_0e[u])=-\omega^2\, u\quad &\mbox{\rm in}\quad \Omega \\
C_0e[u]\nu=0, \quad &\mbox{\rm on}\quad \partial \Omega.
\end{aligned}
\right.
\end{equation}  
Note that the term $C_0e[u]\nu$ can also be written by
\[
C_0e[u]\nu=\lambda_0(\nabla\cdot u) \nu+ 2\mu_0 \sigma(u)\nu.
\]
For the sake of simplicity we shall also denote by $\partial_{\mathbf{\nu}} u:=C_0e[u]\nu $. Recall the following elementary result (cf.\cite{DLL}):
\begin{equation}\label{eq:tract}
\partial_{\mathbf{\nu}}  \mathbf{S}^{\omega}[\mathbf{T}_n^m]|_{-} = (\mathfrak{b}_n-1)\mathbf{T}_n^m, 
\end{equation}
\begin{equation}\label{eq:traci}
  \partial_{\mathbf{\nu}}  \mathbf{S}^{\omega}[ \mathbf{M}_{n}^m] |_{-} =  (\mathfrak{c}_{1n}-1) \mathbf{M}_{n}^m +  \mathfrak{d}_{1n} \tilde{\mathbf{N}}_{n+1}^m,
\end{equation}
and
\begin{equation}\label{eq:tracn}
  \partial_{\mathbf{\nu}}  \mathbf{S}^{\omega}[ \tilde{\mathbf{N}}_{n+1}^m] |_{-} =  \mathfrak{c}_{2n} \mathbf{M}_{n}^m +  (\mathfrak{d}_{2n}-1)\tilde{\mathbf{N}}_{n+1}^m,
\end{equation}
where
\begin{equation}\label{eq:cotrac}
\begin{split}
 \mathfrak{b}_n=&  -i k_s j_{n} (k_s) (k_sh_{n}^{\prime}(k_s) - h_{n}(k_s)), \\
 \mathfrak{c}_{1n}=& -2(n-1)i  \left( \frac{j_{n-1}(k_s) h_{n-1}(k_s) k_s (n+1)}{2n+1}  +  \frac{j_{n-1}(k_p) h_{n-1}(k_p) k_p\mu_0 n}{(\lambda_0+2\mu_0)(2n+1)}   \right)\\
         & + i  \left( \frac{ j_{n-1}(k_s) h_{n}(k_s) k_s^2  (n+1) + j_{n-1}(k_p) h_{n}(k_p) k_p^2 n}{2n+1}   \right), \\
\mathfrak{d}_{1n}=&  2n(n+2) i  \left( \frac{j_{n-1}(k_s) h_{n+1}(k_s) k_s}{2n+1}  -   \frac{j_{n-1}(k_p) h_{n+1}(k_p) k_p\mu_0 }{(\lambda_0+2\mu_0)(2n+1)}  \right)\\
        & +  n i  \left( \frac{ -j_{n-1}(k_s) h_{n}(k_s) k_s^2  + j_{n-1}(k_p) h_{n}(k_p) k_p^2}{2n+1}   \right), \\ 
\mathfrak{c}_{2n}=& - 2(n^2-1)  i  \left( \frac{j_{n+1}(k_s) h_{n-1}(k_s) k_s}{2n+1}  -  \frac{j_{n+1}(k_p) h_{n-1}(k_p) k_p\mu_0 }{(\lambda_0+2\mu_0)(2n+1)}   \right)\\
         & -(n+1) i  \left( \frac{ -j_{n-1}(k_s) h_{n}(k_s) k_s^2   + j_{n-1}(k_p) h_{n}(k_p) k_p^2 }{2n+1}   \right), \\
\mathfrak{d}_{2n}=&  2(n+2)i  \left( \frac{j_{n+1}(k_s) h_{n+1}(k_s) k_s n}{(2n+1)}  +  \frac{j_{n+1}(k_p) h_{n+1}(k_p) k_p \mu_0 (n+1)}{(\lambda_0+2\mu_0)(2n+1)}  \right)\\
         & - i  \left( \frac{ j_{n+1}(k_s) h_{n}(k_s) k_s^2 n  + j_{n+1}(k_p) h_{n}(k_p) k_p^2(n+1) }{2n+1}   \right).
\end{split}
\end{equation}
Similarly, consider the linear combination of  $\mathbf{M}_{n}^m$ and $\mathbf{N}_{n+1}^m$, there should satisfy
\begin{equation}\label{eq:deteig01}
\left|
\begin{array}{cccc}
   \mathfrak{c}_{1n}^r-1 &-\mathfrak{c}_{1n}^i&\mathfrak{d}_{1n}^r  & -\mathfrak{d}_{1n}^i \\
   \mathfrak{c}_{1n}^r-1 &\mathfrak{c}_{1n}^i&\mathfrak{d}_{1n}^r  & \mathfrak{d}_{1n}^i \\
   \mathfrak{c}_{2n}^r &-\mathfrak{c}_{2n}^i&\mathfrak{d}_{2n}^r-1  & -\mathfrak{d}_{2n}^i \\
  \mathfrak{c}_{2n}^r &\mathfrak{c}_{2n}^i&\mathfrak{d}_{2n}^r-1  & \mathfrak{d}_{2n}^i 
\end{array}
\right|=0    
\end{equation}

where $\mathfrak{c}_{jn}^r$, $\mathfrak{d}_{jn}^r$, $j=1,2$ are real parts of $\mathfrak{c}_{jn}$ and $\mathfrak{d}_{jn}$, respectively. where $\mathfrak{c}_{jn}^i$, $\mathfrak{d}_{jn}^i$, $j=1,2$ are imaginary parts of $\mathfrak{c}_{jn}$ and $\mathfrak{d}_{jn}$, respectively. 

Besides, consider $\mathbf{T}_n^m$, then it requires
\[
\mathfrak{b}_{n}-1=0,
\]
which is
\begin{equation}\label{eq:9.15}
\left\{
\begin{aligned}
&k_sj_n'(k_s)=j_n(k_s)\\    
&k_sj_n(k_s)(k_sy_n'(k_s)-y_n(k_s))=1.
\end{aligned}
\right.
\end{equation}
By combing the Wronskian identity
\[
j_n'(t)h_n(t)-j_n(t)h_n'(t)=-\frac{i}{t^2},
\]
one only requires the first equation in (\ref{eq:9.15}) to be satisfied, then the second one satisfies automatically. Thus for such mode, $k_s$ should be the roots of $tj_n'(t)-j_n(t)$ and the associated eigenvalues are $-k_s^2\mu_0^2$.

\section{Asymptotic behavior of Lam\'e operator's eigenvalues {\color{magenta}in charge of Youjun}} 
In this section, we shall consider the asymptotic behavior of the eigenvalues to the Lam\'e operator. Let us focus on the traction free boundary condition.

To begin with, recall the following asymptotic results for spherical Bessel function $j_n(t)$ and Hankel function $h_n(t)$ (cf. \cite{CK}):
\[
j_n(t)=\frac{t^n}{1\cdot 3\cdots (2n+1) }\Big(1+O(\frac{1}{n})\Big), \quad n\rightarrow +\infty
\]
holds in any subset of $\mathbb{R}$ and
\[
h_n(t)=\frac{1\cdot 3\cdots (2n-1)}{it^{n+1}}\Big(1+O(\frac{1}{n})\Big), \quad n\rightarrow +\infty
\]
holds on compact subsets of $(0,+\infty)$. Furthermore, there holds
\[
h_n(t)=\frac{1}{t}e^{i(t-\frac{n+1}{2}\pi)}\Big(1+O(\frac{1}{t})\Big), \quad t\rightarrow +\infty.
\]
Recall also that when $n$ goes to infinity and $t>n+1/2$, then it holds that (see \cite{Kor02}, pp129)
\[
h_n(t)=\sqrt{\frac{1}{t\sqrt{t^2-(n+1/2)^2}}}e^{i\theta_{n}(t)}(1+o(1)),
\]
where
\[
\theta_n(t)=\sqrt{t^2-(n+1/2)^2}-\frac{(n+1)\pi}{2}+(n+1/2)\arcsin{(n+1/2)}/t.
\]

Suppose the trace of eigenfunction on $\partial \Omega$ is $\mathbf{T}_n^m$ (divergence free mode), then the eigenvalues satisfy the following equation:
\begin{equation}\label{eq:eigenLa01}
k_sj_n'(k_s)-j_n(k_s)=0. 
\end{equation}
Suppose $k_s\rightarrow \infty$, then for any fixed $n$, there holds the following relation
\[
\sin(k_s-\frac{n+1}{2}\pi)=O(\frac{1}{k_s}).
\]
Thus the eigenvalues $k_s$ obeys the following asymptotic expansions
\[
k_s=\frac{n+1}{2}\pi+k\pi+O(\frac{1}{k_s}), \quad k\in\mathbb{N}.
\]
Next, suppose $n\rightarrow \infty$, then to ensure the existence of solution to (\ref{eq:eigenLa01}), $k_s$ should also goes to $\infty$. Here we suppose $k_s>n+1/2$, then it holds that
\[
\Theta(k_s,n)\sin{(k_s\Theta(k_s,n)-\frac{n+1}{2}\pi+(n+1/2)\arcsin{\frac{n+1/2}{k_s}})}=O(\frac{1}{k_s^2}),
\]
where
\[
\Theta(k_s,n)=\sqrt{1-(n+1/2)^2/k_s^2}.
\]

Next, suppose the trace of eigenfunction on $\partial \Omega$ is linear combinations of $\mathbf{M}_n^m$ and $\mathbf{N}_{n+1}^m$, then the eigenvalues should satisfy the equation (\ref{eq:deteig01}).
To this end, we always assume $n$ is sufficiently large and we search for the eigenvalues $-\omega^2$ with $\omega$ large enough such that $k_s,k_p>n+3/2$. 

The equation (\ref{eq:deteig01}) can be rewritten by
\begin{equation}\label{eq:eigeneq01}
\Big((\mathfrak{c}_{1n}^r-1)(\mathfrak{d}_{2n}^r-1 )-
\mathfrak{c}_{2n}^r\mathfrak{d}_{1n}^r\Big)
\Big(
\mathfrak{c}_{1n}^i\mathfrak{d}_{2n}^i-\mathfrak{c}_{2n}^i\mathfrak{d}_{1n}^i
\Big)=0.
\end{equation}
Note that
\[
tj_m(t)h_n(t)=\frac{1}{t}(\Theta(t,m)\Theta(t,n))^{-1/2}\\cos{\theta_m(t)}e^{i\theta_n(t)}(1+o(1)).
\]
Suppose $t\gg n$, then one has
\[
tj_m(t)h_n(t)=\frac{1}{t}(\cos(t-\frac{m+1}{2}\pi)e^{i(t-\frac{n+1}{2}\pi)}(1+o(1)).
\]
Thus in the case that $k_s, k_p\gg n$, there holds
\[
(\mathfrak{c}_{1n}^r-1)(\mathfrak{d}_{2n}^r-1 )-
\mathfrak{c}_{2n}^r\mathfrak{d}_{1n}^r=(\frac{1}{2}\cos^4(k_s-\frac{n}{2}\pi)+\frac{1}{2}\cos^4(k_p-\frac{n}{2}\pi)-1)(1+o(1)),
\]
and
\[
\mathfrak{c}_{1n}^i\mathfrak{d}_{2n}^i-\mathfrak{c}_{2n}^i\mathfrak{d}_{1n}^i=(\frac{1}{8}\sin^2(2k_s-n\pi)+\frac{1}{8}\sin^2(2k_p-n\pi))(1+o(1)).
\]
Thus the equation (\ref{eq:eigeneq01}) implies
\[
k_s-k_p=\frac{K}{2}\pi, 
\]
for sufficient large nature number $K\in \mathbb{N}$. In other words,
\[
\omega=\frac{K\pi}{1/\sqrt{\mu_0}-1/\sqrt{\lambda_0+2\mu_0}}, \quad K\in \mathbb{N}.
\]

\end{comment}

\section{Special eigenfunctions for the instantaneous
term}
\setcounter{equation}{0}
From this section, we prepare to generate the C-ev's for the FOE. Hence, in this section, we assume that the Earth $\Omega$ is an open ball with radius $R>0$ centered at the origin, and consider it as a homogeneous isotropic EBM. To begin with, we look for special radial functions that are eigenfunctions of the volumetric/deviatoric parts of the instantaneous part simultaneously, and satisfy the traction-free boundary condition for the instantaneous part. 

Let $Q_A$ and $Q_B$ be the operators defined as
\begin{equation}\label{s3.1}
\begin{aligned}
Q_Au:=(\lambda_0+\frac{2}{3}\mu_0)\nabla\cdot \big((\nabla\cdot u) I_3\big)=(\lambda_0+\frac{2}{3}\mu_0)\nabla (\nabla \cdot u)
\end{aligned}
\end{equation}
and
\begin{equation}\label{s3.2}
\begin{aligned}
Q_Bu:=2\mu_0\nabla\cdot \big(e[u]-3^{-1}(\nabla\cdot u) I_3\big).
\end{aligned}
\end{equation}
Also, let $H(p)$ be the Hessian of $p$ given as $H(p)=(\partial^2_{x_i,x_j}p)_{i,j=1}^3$. 
\begin{lemma}\label{s_lem3.1}
Let $u$ satisfy
\begin{equation}\label{s3.3}
\begin{array}{l}
\begin{cases}
Q_Au=-(\lambda_0+\frac{2}{3}\mu_0)k_a\, u\quad &\mbox{\rm in}\quad \Omega,\\
Q_Bu=-\frac{4}{3}\mu_0k_a\, u\quad &\mbox{\rm in}\quad \Omega,\\
\big(\lambda_0(\nabla \cdot u) I_3+2\mu_0e[u]\big)\nu=0      & {\rm on}\quad \partial\Omega
\end{cases}
\end{array} 
\end{equation}
for some constant $k_a>0$.
Then, $p:=\nabla\cdot u$ satisfies
\begin{equation}\label{s3.4}
\begin{array}{l}
\begin{cases}
\Delta p=-k_a\, p&\mbox{\rm in}\quad \Omega,\\
\big(\lambda_0p\tilde I-2\mu_0k_a^{-1}H(p)\big)\nu=0    & {\rm on}\quad \partial\Omega.
\end{cases}
\end{array}
\end{equation}
\end{lemma}
\begin{proof}
Direct computations give that
\begin{equation}\label{s3.5}
\begin{aligned}
 (\lambda_0+\frac{2}{3}\mu_0)\Delta p=\nabla \cdot (Q_A u)=-(\lambda_0+\frac{2}{3}\mu_0)k_a\, p
\end{aligned}
\end{equation}
and
\begin{equation}\label{s3.6}
\begin{aligned}
 2\mu_0(1-3^{-1})\Delta p=\nabla \cdot (Q_B(x)u)= -\frac{4}{3}\mu_0k_a\, p.
\end{aligned}
\end{equation}
Adding \eqref{s3.5} and \eqref{s3.6}, we have the first equation of \eqref{s3.4}. Using the first equation of \eqref{s3.3}, we have  
$$\nabla p=\nabla (\nabla \cdot u)=-k_au,\,\,\text{and hence}\,\,e[u]=-k_a^{-1}H(p).$$
Then, plugging these into the third equation of \eqref{s3.3}, we have the second equation of \eqref{s3.4}.
\end{proof}
  
\medskip\noindent
On the other hand, the converse of Lemma \ref{s_lem3.1} is also true. More precisely, we have the following lemma.

\begin{lemma}\label{s_lem3.2}
Let  $p$ satisfy
\begin{equation}\label{s3.7}
\begin{array}{l}
\begin{cases}
\Delta p=-k_b\, p&\mbox{\rm in}\quad \Omega,\\
\big(\lambda_0p I_3 -2\mu_0k_b^{-1}H(p)\big)\nu=0    & {\rm on}\quad \partial\Omega,
\end{cases}
\end{array}
\end{equation}
where $k_b>0$ is a constant. 
Then, $u:=-k_b^{-1}\nabla p$ satisfies
\begin{equation}\label{s3.8}
\begin{array}{l}
\begin{cases}
Q_Au=-(\lambda_0+\frac{2}{3}\mu_0)k_b\, u\quad &\mbox{\rm in}\quad \Omega,\\
Q_Bu=-\frac{4}{3}\mu_0k_b\, u\quad &\mbox{\rm in}\quad \Omega,\\
\big(\lambda_0(\nabla \cdot u) I_3+2\mu_0e[u]\big)\nu=0     & {\rm on}\quad \partial\Omega.
\end{cases}
\end{array} 
\end{equation}
Further, we have $p=\nabla\cdot u$, which implies the equivalence of \eqref{s3.3} and \eqref{s3.4}.
\end{lemma}

\begin{proof}
Direct computations give that
\begin{equation}\label{s3.9}
\begin{array}{ll}
Q_A(x)u&=(\lambda_0+\frac{2}{3}\mu_0)\nabla (\nabla\cdot u) 
=-(\lambda_0+\frac{2}{3}\mu_0)k_b^{-1}\nabla \Delta p\\
&=(\lambda_0+\frac{2}{3}\mu_0)\nabla p
=-(\lambda_0+\frac{2}{3}\mu_0)k_b\,u
\end{array}
\end{equation}
and
\begin{equation}\label{s3.10}
\begin{array}{ll}
Q_B(x)u&=\mu_0\nabla (\nabla\cdot u)+\mu_0 \Delta u-\frac{2}{3}\mu_0 \nabla (\nabla\cdot u)\\
&=-\mu_0k_b^{-1}\nabla \Delta p-\mu_0 k_b^{-1}\nabla \Delta p+\frac{2}{3}\mu_0 k_b^{-1}\nabla \Delta p\\
&=\frac{4}{3}\mu_0\nabla p
=-\frac{4}{3}\mu_0k_b\,u.
\end{array}
\end{equation}
Also, we have
$$-k_b\nabla \cdot u=\nabla \cdot \nabla p=\Delta p=-k_b p$$
which implies $p=\nabla \cdot u$. Then, we can finish the proof by showing the third equation of \eqref{s3.8} by direct computations.
\end{proof}

\begin{comment}
If $u$ is a solution to \eqref{s3.3}, then from Lemma \ref{s_lem3.1}, $p=\nabla \cdot u$ is a solution to \eqref{s3.4}. Actually, from the first equation of \eqref{s3.3}, we have $\nabla (\nabla \cdot u)=-k_au$.  From Lemma \ref{s_lem3.2}, $\nabla (\nabla \cdot u)$ is a solution to \eqref{s3.8}. 
Thus, we know that \eqref{s3.3} and \eqref{s3.7} are equivalent. 
\end{comment}

Now, we seek $p$ as a radial function $p(x)=p_1(r)$ with $r=|x|$ to satisfy \eqref{s3.7}. By direct computations, we have
$$H(p)\nu=\nabla p_1'(r)=p_1''(r)\nu.$$
Thus, \eqref{s3.7} is transformed to 
\begin{equation}\label{s3.11}
\begin{array}{l}
\begin{cases}
r^2 p_1''(r)+2rp_1'(r)=-k_br^2 p_1(r)&\mbox{\rm for}\quad 0\leq r< R,\\
\lambda_0p_1(r)-2\mu_0k_b^{-1}p_1''(r)=0    & {\rm for}\quad r=R.
\end{cases}
\end{array}
\end{equation}
Using the first equation of \eqref{s3.11}, we have
\begin{equation}\label{s3.12}
\begin{aligned}
\lambda_0p_1(r)-2\mu_0k_b^{-1}p_1''(r)&=\lambda_0p_1(r)-2\mu_0k_b^{-1}(-k_bp_1(r)-2r^{-1}p_1'(r))\\
&=(\lambda_0+2\mu_0)p_1(r)+4\mu_0k_b^{-1}r^{-1}p_1'(r)=0.
\end{aligned}
\end{equation}
Hence, using \eqref{s3.12}, \eqref{s3.11} becomes
\begin{equation}\label{s3.13}
\begin{array}{l}
\begin{cases}
r^2 p_1''(r)+2rp_1'(r)+k_br^2 p_1(r)=0,&\mbox{\rm for}\quad 0\leq r< R,\\
(\lambda_0+2\mu_0)p_1(r)+4\mu_0k_b^{-1}r^{-1}p_1'(r)=0    & {\rm for}\quad r=R.
\end{cases}
\end{array}
\end{equation}
By putting $\sqrt{k_b} r=\eta$ and $p_2(\eta)=p_1(\eta/\sqrt{k_b})$, we rewrite \eqref{s3.13} as
\begin{equation}\label{s3.14}
\begin{array}{l}
\begin{cases}
\eta^2 p_2''(\eta)+2\eta p_2'(\eta)+\eta^2 p_2(\eta)=0,&\mbox{\rm for}\quad 0\leq \eta< \sqrt{k_b}R,\\
(\lambda_0+2\mu_0)p_2(\eta)+4\mu_0\eta^{-1}p_2'(\eta)=0    & {\rm for}\quad \eta=\sqrt{k_b}R.
\end{cases}
\end{array}
\end{equation}

Next, consider the spherical Bessel function $J(\eta)$ of order $0$ given as
\begin{equation}\label{s3.15}
\begin{array}{l}
J(\eta):=\frac{\sin \eta}{\eta}=\displaystyle\sum_{j=0}^\infty\frac{(-1)^j}{(2j+1)!}\eta^{2j},
\end{array}
\end{equation}
which is analytic in ${\mathbb R}$ (see \cite{Kor02}). Then, $J(\eta)$ satisfies
\begin{equation}\label{s3.16}
\begin{aligned}
\eta^2 J''(\eta)+2\eta J'(\eta)+\eta^2 J(\eta)=0\,\,\,\text{for}\,\,\,\eta\in{\mathbb R}.
\end{aligned}
\end{equation}
Comparing the first equation of \eqref{s3.14} and \eqref{s3.16}, we consider $p_2(\eta)$ as $J(\eta)$.

Next, concerning the second equation of \eqref{s3.14}, we search for positive roots of  $$(\lambda_0+2\mu_0)J(\eta)+4\mu_0\eta^{-1}J'(\eta)=0,$$
which is equivalent to search for positive roots of the following equation
\begin{equation}\label{s3.17}
\begin{aligned}
f(\eta):
=&\big((\lambda_0+2\mu_0)\eta^2-4\mu_0\big)\sin\eta+4\mu_0\eta\cos\eta=0.
\end{aligned}
\end{equation}

\medskip\noindent
Then, we have the following lemma.
\begin{lemma}\label{s_lem3.3}
There exists a unique $r_\ell\in((\ell-1)\pi, \ell \pi)$ such that $f(r_\ell)=0$ for each  $\ell=1,2,\cdots$. In fact, we have
$r_l\in((l-\frac{1}{2})\pi, l\pi)$ for each  $l=1,2,\cdots$. 
\end{lemma}
\begin{proof}
By direct computations, we have
\begin{equation}\label{s3.18}
\begin{aligned}
f'(\eta)=(\lambda_0+2\mu_0)\eta^2\cos\eta+2\lambda_0\eta\sin\eta
\end{aligned}
\end{equation}
and
\begin{equation}\label{s3.19}
\begin{aligned}
f''(\eta)=4(\lambda_0+\mu_0)\cos\eta+\big(2\lambda_0-(\lambda_0+2\mu_0)\eta^2\big)\sin\eta.
\end{aligned}
\end{equation}
We divide $\ell$ into two cases.

\medskip
\noindent
Case 1: $\ell=2m-1$ with a positive integer $m$.

From \eqref{s3.18}, we have that 
\begin{equation}\label{s3.20}
\begin{array}{l}
\begin{cases}
f'(\eta)>0,&\mbox{\rm if}\quad \eta\in\big((2m-2)\pi, (2m-\frac{3}{2})\pi\big),\\
f'(\eta)<0    & {\rm if}\quad \eta=(2m-1)\pi.
\end{cases}
\end{array}
\end{equation}
Since
\begin{equation*}
\begin{aligned}
2\lambda_0-(\lambda_0+2\mu_0)\eta^2<0\,\,\text{for}\,\,\eta\in[(2m-\frac{3}{2})\pi,(2m-1)\pi]
\end{aligned}
\end{equation*}
if $\eta\in[(2m-\frac{3}{2})\pi,(2m-1)\pi]$. Thus, we obtain that
\begin{equation}\label{s3.21}
\begin{aligned}
f''(\eta)<0
\end{aligned}
\end{equation}
if $\eta\in\big((2m-\frac{3}{2})\pi,(2m-1)\pi)\big)$. 
By \eqref{s3.20} and \eqref{s3.21}, there exists a unique $\eta_{2m-1}\in \big((2m-\frac{3}{2})\pi,(2m-1)\pi\big)$ such that $f'(\eta_{2m-1})=0$. We also have that
\begin{equation}\label{s3.22}
\begin{array}{l}
\begin{cases}
f'(\eta)>0,&\mbox{\rm if}\quad \eta\in((2m-2)\pi, \eta_{2m-1}),\\
f'(\eta)<0    & {\rm if}\quad \eta\in(\eta_{2m-1}, (2m-1)\pi).
\end{cases}
\end{array}
\end{equation}
Hence, by $f((2m-2)\pi)> 0$, $f((2m-1)\pi)< 0$ and \eqref{s3.22}, there exists a unique $r_{2m-1}\in (\eta_{2m-1},(2m-1)\pi)$ such that $f(r_{2m-1})=0$.

\medskip\noindent
Case 2: $\ell=2m$ with a positive integer $m$.

From \eqref{s3.18}, we have 
\begin{equation}\label{s3.23}
\begin{array}{l}
\begin{cases}
f'(\eta)<0,&\mbox{\rm if}\quad \eta\in\big((2m-1)\pi, (2m-\frac{1}{2})\pi\big),\\
f'(\eta)>0    & {\rm if}\quad \eta=2m\pi.
\end{cases}
\end{array}
\end{equation}
Since
\begin{equation*}
\begin{aligned}
2\lambda_0-(\lambda_0+2\mu_0)\eta^2<0\,\,\,\,\text{for}\,\,\,\,
\eta\in[(2m-\frac{1}{2})\pi,2m\pi],
\end{aligned}
\end{equation*}
we have
\begin{equation}\label{s3.24}
\begin{aligned}
f''(\eta)>0\,\,\,\,\text{for}\,\,\,\,\eta\in[(2m-\frac{1}{2})\pi,2m\pi)].
\end{aligned}
\end{equation}
By \eqref{s3.23} and \eqref{s3.24}, there exists a unique $\eta_{2m}\in \big((2m-\frac{1}{2})\pi,2m\pi\big)$ such that $f'(\eta_{2m})=0$. We also have
\begin{equation}\label{s3.25}
\begin{array}{l}
\begin{cases}
f'(\eta)<0,&\mbox{\rm if}\quad \eta\in((2m-1)\pi, \eta_{2m}),\\
f'(\eta)>0    & {\rm if}\quad \eta\in(\eta_{2m}, 2m\pi).
\end{cases}
\end{array}
\end{equation}
Hence, by $f((2m-1)\pi)< 0$, $f(2m\pi)> 0$ and \eqref{s3.25}, there exists a unique $r_{2m}\in (\eta_{2m},2m\pi)$ such that $f(r_{2m})=0$.

\end{proof}

\medskip\noindent
Based on the previous lemma, the following lemma states that we can completely solve \eqref{s3.13} by taking $k_b=\frac{r_\ell^2}{R^2}$ for  $\ell=1,2,\cdots$.
\begin{lemma}\label{s_lem3.4}
For $\ell=1,2,\cdots$, let $y_\ell(r)=J(\frac{r_\ell}{R}r)$. Then $y_\ell$ satisfies
\begin{equation}\label{s3.26}
\begin{array}{l}
\begin{cases}
r^2 y_\ell''(r)+2ry_\ell'(r)=-k_br^2 y_\ell(r)&\mbox{\rm for}\quad 0\leq r< R,\\
(\lambda_0+2\mu_0)y_\ell(r)+4\mu_0k_b^{-1}r^{-1}y_\ell'(r)=0     & {\rm for}\quad r=R.
\end{cases}
\end{array}
\end{equation}
\end{lemma}
\begin{proof}
The first equation of \eqref{s3.26} follows from \eqref{s3.16}. As for the second equation of \eqref{s3.26}, it follows from $f(r_\ell)=0$ for $\ell=1,2,\cdots$.
\end{proof}
Thus, we have solved \eqref{s3.13} to get $p(x)=p_1(r)$. Then, we can get $u=-k_b^{-1} \nabla p$   
satisfying \eqref{s3.8}. 
Also, by direct computations, we have that
\begin{equation}\label{s3.27}
\begin{array}{ll}
u(x)&=-k_b^{-1}  y'_\ell(r) \frac{x}{r}=-\frac{R^2}{r_\ell^2}\big(r\cos(\frac{r_\ell}{R}r)-\frac{R}{r_\ell}\sin (\frac{r_\ell}{R}r)\big)\frac{x}{r^3}\\
&=-l_b^{-3}\{\displaystyle\sum_{i=0}^\infty\big((2i)!\big)^{-1}(-1)^i r(l_b r)^{2i+1}+\displaystyle\sum_{j=1}^\infty\big((2j-1)!\big)^{-1} (-1)^j(l_br)^{2j-1}\}\frac{x}{r^3}
\end{array}
\end{equation}
\begin{comment}    
\begin{equation}\label{s3.27}
u(x)=
\begin{array}{l}
\begin{cases}
-k_b^{-1}  y'_l(r) \frac{x}{r}=-\frac{R^2}{r_l^2}\big(r\cos(\frac{r_l}{R}r)-\frac{R}{r_l}\sin (\frac{r_l}{R}r)\big)\frac{x}{r^3}&\mbox{\rm if}\quad x\neq 0,\\
0    & {\rm if}\quad x=0,
\end{cases}
\end{array}
\end{equation}
\end{comment}
for $x\in\overline\Omega$, where $r=|x|$, $l_b=\sqrt{k_b}$ and $r_\ell$, $\ell=1,2,\cdots$ satisfy 
\begin{equation}\label{s3.28}
\begin{aligned}
f(r_\ell)=(\lambda_0+2\mu_0)r^2_\ell\sin r_\ell+4\mu_0(r_\ell\cos r_\ell-\sin r_\ell)=0.
\end{aligned}
\end{equation}
Note here that, since the first terms of the summations in \eqref{s3.27} cancel out, $u(x)$ is analytic on $\overline\Omega$. Therefore, $u(x)$ of \eqref{s3.27} satisfies
\begin{equation}\label{s3.29}
\begin{array}{l}
\begin{cases}
Q_Au=-(\lambda_0+\frac{2}{3}\mu_0)\frac{r_\ell^2}{R^2}\, u\quad &\mbox{\rm in}\quad \Omega,\\
Q_Bu=-\frac{4}{3}\mu_0\frac{r_\ell^2}{R^2}\, u\quad &\mbox{\rm in}\quad \Omega,\\
\big(\lambda_0(\nabla \cdot u)\tilde I+2\mu_0e[u]\big)\nu=0     & {\rm on}\quad \partial \Omega
\end{cases}
\end{array} 
\end{equation}
for $\ell=1,2\cdots$.

\begin{comment}
\begin{lemma}\label{s_lem3.5}
$\partial_{x_j}u_i(x)$ are continuous in $x\in \Omega$ for $1\leq i,j\leq 3$.
\end{lemma}
\begin{proof}
From \eqref{s3.27} and the first equation of \eqref{s3.26}, we have for $r\neq 0$ that
\begin{equation}\label{s3.30}
\begin{aligned}
\partial_{x_j}u_i(x)&=-k_b^{-1}  y''_l(r) \frac{x_ix_j}{r^2}-k_b^{-1}  y'_l(r)( -\frac{x_ix_j}{r^3}+\delta_{ij}\frac{1}{r})\\
&=-k_b^{-1}  \big(y''_l(r)-\frac{y'_l(r)}{r}\big)\frac{x_ix_j}{r^2}-k_b^{-1}\frac{y'_l(r)}{r}\delta_{ij}\\
&=-k_b^{-1}  \big(-\frac{3y'_l(r)}{r}-k_by_l(r)\big)\frac{x_ix_j}{r^2}-k_b^{-1}\frac{y'_l(r)}{r}\delta_{ij}.
\end{aligned}
\end{equation}
From the definition of $\partial_{x_j}u_i(0)$, we have that
\begin{equation}\label{s3.31}
\begin{aligned}
\partial_{x_j}u_i(0)=\lim_{r\rightarrow 0}-k_b^{-1}\frac{y'_l(r)}{r}\delta_{ij}=\frac{1}{3}\delta_{ij}.
\end{aligned}
\end{equation}
The direct computations give that
\begin{equation}\label{s3.32}
\begin{aligned}
\lim_{r\rightarrow 0} \big(-\frac{3y'_l(r)}{r}-k_by_l(r)\big)=k_b-k_b=0.
\end{aligned}
\end{equation}
Combining \eqref{s3.30}, \eqref{s3.31} and \eqref{s3.32}, we obtain that
\begin{equation}\label{s3.33}
\begin{aligned}
\lim_{x\rightarrow 0}\partial_{x_j}u_i(x)=\partial_{x_j}u_i(0)=\frac{1}{3}\delta_{ij}
\end{aligned}
\end{equation}
which implies $\partial_{x_j}u_i(x)$ are continuous in $x\in \Omega$ for $1\leq i,j\leq 3$.

\end{proof}
\end{comment}

\begin{comment}
Moreover, we can solve $p(x)$ by solving the following equations.
\begin{equation}\label{I8.23}
\begin{array}{l}
\begin{cases}
p(x)=p_1(r), & (1-\delta) R \leq r=|x| \leq R,\\
\Delta p(x)=k_b\, p(x),&\mbox{\rm in}\quad  |x|< (1-\delta)R\\
p(x)=p_1((1-\delta)R)    & {\rm on}\quad |x|=(1-\delta)R,
\end{cases}
\end{array}
\end{equation}
here $\delta$ is a constant satisfying $0<\delta\leq 1$.

\eqref{I8.14} means that we can find a solution $p$ which behaves like the solution of the Helmholtz equation with constant boundary conditions.
\end{comment}

\begin{comment}
By considering different boundary conditions of \eqref{s3.3}, we can solve the associated problems of \eqref{s3.7}. 

If
\begin{equation}\label{s3.30}
\begin{array}{l}
\begin{cases}
Q_Au=(\lambda_0+\frac{2}{3}\mu_0)k_a\, u\quad &\mbox{\rm in}\quad \Omega,\\
Q_Bu=\frac{4}{3}\mu_0k_a\, u\quad &\mbox{\rm in}\quad \Omega,\\
(\nabla \cdot u)\tilde I\nu=0      & {\rm on}\quad \partial\Omega
\end{cases}
\end{array} 
\end{equation}
then we consider
\begin{equation}\label{s3.31}
\begin{array}{l}
\begin{cases}
\Delta p(x)=k_a\, p(x),&\mbox{\rm in}\quad \Omega,\\
p=0    & {\rm on}\quad \partial\Omega.
\end{cases}
\end{array}
\end{equation}
On the other hand, if
\begin{equation}\label{s3.32}
\begin{array}{l}
\begin{cases}
Q_Au=(\lambda_0+\frac{2}{3}\mu_0)k_a\, u\quad &\mbox{\rm in}\quad \Omega,\\
Q_Bu=\frac{4}{3}\mu_0k_a\, u\quad &\mbox{\rm in}\quad \Omega,\\
u \cdot \nu=0      & {\rm on}\quad \partial\Omega
\end{cases}
\end{array} 
\end{equation}
then we consider
\begin{equation}\label{s3.33}
\begin{array}{l}
\begin{cases}
\Delta p(x)=k_a\, p(x),&\mbox{\rm in}\quad \Omega,\\
\nabla p\cdot \nu=0    & {\rm on}\quad \partial\Omega.
\end{cases}
\end{array}
\end{equation}
\end{comment}






\section{C-ev's for the FOE }
\setcounter{equation}{0}

In this section, we consider cluster eigenvalues for the FOE using the special eigenfunctions for the instantaneous part given in the previous section. First of all, from \eqref{s2.28}, \eqref{s3.1} and \eqref{s3.2}, the BVS becomes
\begin{equation}\label{s4.1}
\begin{array}{rcl}
\partial_t^2 u(t) &=&Q_Au(t)+Q_Bu(t)-\displaystyle\sum_{j=0}^n \kappa_j(q_0^j)^2\int_0^t e^{-\kappa_j(t-s)}Q_Au(s)\,ds\\
&&-\displaystyle\sum_{j=0}^n \tau_j(v_0^j)^2\int_0^t e^{-\tau_j(t-s)}Q_Bu(s)\,ds\qquad\text{in $(0,\infty)\times\Omega$}.
\end{array}
\end{equation}
Here and hereafter, note that likewise before, we suppressed the variable $x$.
Without loss of the generality, we assume that
\begin{equation}\label{s4.2}
\begin{aligned}
\tau_0<\tau_1<\cdots<\tau_n<\kappa_0<\kappa_1<\cdots<\kappa_n.
\end{aligned}
\end{equation}
For simplicity, we introduce the following shorthand notation 
\begin{equation}\label{s4.3}
\left\{
\begin{array}{ll}\beta_{j+1}=\tau_j,\, \alpha_{j+1}=\frac{4}{3}\mu_0\tau_j(v_0^j)^2,\,  w_{j+1}:=(\frac{4}{3}\mu_0)^{-1}\int_0^t e^{-\tau_j(t-s)}Q_Bu(s)\,ds,\,\beta_{j+n+2}=\kappa_j,
 \\
\\
 \alpha_{j+n+2}=(\lambda_0+\frac{2}{3}\mu_0)\kappa_j(q_0^j)^2, \,\,w_{j+n+2}:=(\lambda_0+\frac{2}{3}\mu_0)^{-1}\int_0^t e^{-\kappa_j(t-s)}Q_Au(s)\,ds
\end{array}
\right.
\end{equation}  
for $0\leq j\leq n$. Then, \eqref{s4.1} takes the form
\begin{equation}\label{s4.4}
\begin{aligned}
\partial_t^2 u=&Q_A u +Q_B u-  \sum_{j=1}^{2n+2} \alpha_jw_j\\
=&Q_A u +Q_B u- \sum_{j=1}^{n+1} \alpha_j\int_0^t e^{-\beta_j(t-s)}\nabla\cdot \big(\frac{3}{2}e[u]-\frac{1}{2}(\nabla\cdot u) I_3\big)(s)\,ds\\
&-\sum_{j=n+2}^{2n+2} \alpha_j\int_0^t e^{-\beta_j(t-s)}\nabla\cdot \big((\nabla\cdot u) I_3\big)(s)\,ds.
\end{aligned}
\end{equation}
Here, we note a little ahead of time that the inverse problem in the next section, in which C-ev's are used as spectral data, is to recover the coefficient $(\lambda_0+2\mu_0)k_b$ of $Q_A u+Q_B u$ with $u$ satisfying \eqref{s3.8} obtained at the end of Section 3, and $\alpha_j$'s, $\beta_j$'s.

Now, by introducing a new dependent variable $v=u'$, we write \eqref{s4.1} in the form of a system of partial differential equations that are of first order in time given as
\begin{equation}\label{s4.5}
\left\{
\begin{array}{rcl}
u'&=&v,\\
v'&=&Q_A u +Q_B u-   \sum_{i=1}^{n+2} \alpha_iw_i,\\
w_j'&=&(\frac{4}{3}\mu_0)^{-1}Q_B u -\beta_jw_j\,\,\,\text{for}\,\,\, 1\leq j \leq n+1\\
w_j'&=&(\lambda_0+\frac{2}{3}\mu_0)^{-1}Q_A u -\beta_jw_j\,\,\,\text{for}\,\,\, n+2\leq j \leq 2n+2.
\end{array}
\right.
\end{equation}
We refer to this system as an augmented system. The matrix-vector representation of this system is \begin{equation}\label{s4.6}
\begin{aligned}
U'=A_{2n+4}U,
\end{aligned}
\end{equation}
where $A_{2n+4}$ is a $2n+4$ square matrix of the second-order partial differential operators in $x$ and $U=(u,v,w_1,\cdots,w_{2n+2})^{\mathfrak{t}}$. More precisely, we consider $A_{2n+4}$ as a densely defined closed operator on the $2n+4$ number of products of $L^2(\Omega)$, denoted by $L^2(\Omega)^{2n+4}$, with the domain
\begin{equation}\label{s4.7}
D(A_{2n+4}):=\left\{
\begin{array}{ll}
U=(u,v,w_1,\cdots,w_N)^{\mathfrak{t}}\in H^2(\Omega)\times H^1(\Omega)\times L^2(\Omega)^{2n+2}:\\
\qquad\big(\lambda_0(\nabla \cdot u)\tilde I+2\mu_0e[u]\big)\nu=0  \quad    {\rm on}\quad \partial \Omega,\,t>0
\end{array}
\right\},
\end{equation}
where $H^s(\Omega),\,s=1,2$ are the $L^2(\Omega)$-based  Sobolev spaces of order $s=1,2$. Then, we have an eigenvalue problem given as
\begin{equation}\label{s4.8}
\begin{aligned}
z U=A_{2n+4}\,U.
\end{aligned}
\end{equation}
The cluster eigenvalues problems associated to \eqref{s4.8} that we are interested in are as follows. Plugging the solution $u$ depending on $\ell$ of \eqref{s3.29} into \eqref{s4.8}, we consider the following reduced eigenvalues problem given as
\begin{equation}\label{s4.9}
\begin{aligned}
z\,U^\ell=A_{2n+4}^\ell\,U^\ell.
\end{aligned}
\end{equation}
More precisely, for each $\ell\in{\mathbb N}$, we look for an eigenvalue $z\not=-\beta_j,\,1\le j\le 2n+2$ of \eqref{s4.9} with the associated eigenfunction $U^\ell=(u,v,w_1,\cdots,w_{2n+2})^{\mathfrak{t}}$ of \eqref{s4.8}. A detail form of \eqref{s4.9} is given as
\begin{equation}\label{s4.10}
\left\{
\begin{array}{ll}
zu=v,\\
zv=-(\lambda_0+2\mu_0)\frac{r_\ell^2}{R^2} u-  \sum_{i=1}^{n+2} \alpha_iw_i,\\
zw_j=-\frac{r_\ell^2}{R^2} u -\beta_jw_j, \qquad\quad&\ 1\leq j \leq 2n+2
\end{array}
\right.
\end{equation}
and the form of matrix $A^\ell_{2n+4}$ can be read from \eqref{s4.8}. It is important to note here that $u$ satisfying \eqref{s4.10} is an eigenfunction of the q-EBM with eigenvalue $z^2$ satisfying the traction-free boundary condition for the q-EBM. This is because $u$ satisfies the instantaneous elastic equation with a traction-free boundary condition and this boundary condition is invariant under adding any zero-order terms with real coefficients. 

From \eqref{s4.10},  we have the following
\begin{equation}\label{s4.11}
\left\{
\begin{array}{ll}
z^2u+(\lambda_0+2\mu_0)\frac{r_\ell^2}{R^2} u+  \sum_{i=1}^{n+2} \alpha_iw_i=0,\\
(z+\beta_j)w_j=-\frac{r_\ell^2}{R^2} u, \qquad\quad&\ 1\leq j \leq 2n+2.
\end{array}
\right.
\end{equation}
Multiplying $\Pi_{1\leq j\leq 2n+2}(z+\beta_j)$ to the first equation of \eqref{s4.11} and using the second and third equations of \eqref{s4.11}, we obtain
\begin{equation}\label{s4.12}
\begin{aligned}
(z^2+(\lambda_0+2\mu_0)\frac{r_\ell^2}{R^2})\Pi_{1\leq j\leq 2n+2}(z+\beta_j) - \frac{r_\ell^2}{R^2} \sum_{i=1}^{2n+2} \alpha_i\Pi_{1\leq j\leq 2n+2,j\neq i}(z+\beta_j)=0,
\end{aligned}
\end{equation}
which is nothing but
\begin{equation}\label{s4.13}
\begin{aligned}
((\lambda_0+2\mu_0)+\frac{R^2}{r_l^2}z^2)\Pi_{1\leq j\leq 2n+2}(z+\beta_j) -  \sum_{i=1}^{2n+2} \alpha_i\Pi_{1\leq j\leq 2n+2,j\neq i}(z+\beta_j)=0.
\end{aligned}
\end{equation}
From the arguments deriving \eqref{s4.12}, it is clear that for each $\ell\in{\mathbb N}$, the roots of \eqref{s4.12} are the eigenvalues of $A_{2n+4}^\ell$ and the associated eigenvectors can be obtained from \eqref{s3.27} and \eqref{s4.5}.
We define the characteristic polynomial $P^\ell_{2n+4}(z)$ associated to $\partial_t-A_{2n+4}^\ell$ and its limit polynomial $P_{2n+2}(z)$ as $\ell\rightarrow\infty$ by
\begin{equation}\label{s4.14}
\begin{aligned}
P^\ell_{2n+4}(z):=&((\lambda_0+2\mu_0)+\frac{R^2}{r_l^2}z^2)\Pi_{1\leq j\leq 2n+2}(z+\beta_j) -  \sum_{i=1}^{2n+2} \alpha_i\Pi_{1\leq j\leq 2n+2,j\neq i}(z+\beta_j)
\end{aligned}
\end{equation}
and
\begin{equation}\label{s4.15}
\begin{aligned}
P_{2n+2}(z):=&(\lambda_0+2\mu_0)\Pi_{1\leq j\leq 2n+2}(z+\beta_j) - \sum_{i=1}^{2n+2} \alpha_i\Pi_{1\leq j\leq 2n+2,j\neq i}(z+\beta_j)
\end{aligned}
\end{equation}
respectively. We note that $P_{2n+2}(z)$ is the charactersite polynomial associated to  
\begin{equation}\label{s4.16}
\left\{
\begin{array}{rcl}
0&=&Q_A u +Q_B u-  \sum_{j=1}^{2n+2} \alpha_jw_j,\\
w_j'&=&(\frac{4}{3}\mu_0)^{-1}Q_B u -\beta_jw_j\,\,\,\text{for}\,\,\, 1\leq j \leq n+1\\
w_j'&=&(\lambda_0+\frac{2}{3}\mu_0)^{-1}Q_A u -\beta_jw_j\,\,\,\text{for}\,\,\, n+2\leq j \leq 2n+2.
\end{array}
\right.
\end{equation}
This can be easily shown in the same way as we derived $P^\ell_{2n+4}(z)$. 






\section{Structure of C-ev's and recovering q-EBM from C-ev's}
\setcounter{equation}{0}
%{\color{red} Let's restrict the structure of C-ev's to what we need for recovering the relaxation tensor from the C-ev's. More precisely, by pointing out the similarity or the correspondence between this paper and Applied Mathematical Letter (AML), we state Lemma in the AML using the notation of this paper. Then, give the formulation of the inverse cluster eigenvalue problem for this paper, and the statement of the inversion corresponding to Theorem 1 in AML.}

In \cite{CDDLN1} and \cite{CDDLN2}, the roots of the characteristic polynomial equation \eqref{s5.1} given below were used to recover the relaxation tensor of the EBM describing its glassy state by the Prony series. Usually, the stretched exponential function is used instead of its approximation, the Prony series (see \cite{CDDLN1, CDDLN2} and references there in).
\begin{equation}\label{s5.1}
\begin{aligned}
P_N^k(\lambda):=(D+\frac{\lambda^2}{(2k-1)^2})\Pi_{1\leq j\leq N}(\lambda+r_j) -  \sum_{i=1}^N b_i\Pi_{1\leq j\leq N,j\neq i}(\lambda+r_j) =0.
\end{aligned}
\end{equation}
Comparing \eqref{s4.14} with \eqref{s5.1}, the notational correspondence between these two characteristic equations is given as
$$\lambda=z,\,N=2n+2,\,D=\lambda_0+2\mu_0,\, (2k-1)^2=\frac{r_\ell^2}{R^2},\,r_j=\beta_j,\,b_j=\alpha_j.$$
Here, the notation $r_\ell$ and $r_j$ is a bit confusing, but they can be distinguished by being aware that the left-hand side and right-hand side of each equation come from \eqref{s5.1} and what we have introduced before, respectively. 
Hence, we can literally transform Lemma 3.2 of \cite{CDDLN1} on the properties of the roots of \eqref{s5.1} to those of \eqref{s4.14}. Thus, we have the following lemma on the structure of C-ev's.
\begin{lemma}\label{s_lem5.1}
There exists at least $2n+2$ real roots $a^\ell_j$, $1\leq j\leq 2n+2$ of  $P^\ell_{2n+4}(z)=0$ such that 
\begin{equation}\label{s5.2}
\begin{aligned}
-\beta_{2n+2}<a^\ell_{2n+2}<-\beta_{2n+1}<a^\ell_{2n+1}<\cdots<-\beta_{2}<a^\ell_2<-\beta_{1}<a^\ell_1
\end{aligned}
\end{equation}
and
\begin{equation}
\label{s5.3}
 \sum_{i=1}^{2n+2} \frac{\alpha_i}{a^\ell_j+\beta_i} =(\lambda_0+2\mu_0)+\frac{R^2}{r_l^2}(a^\ell_j)^2,\,\,1\le j\le 2n+2.
\end{equation}
The other two roots are contained in the set $B^D_-\cup \{(-\beta_{2n+2},a^\ell_1)\} $, where $B^D_-:=\{c+id :\frac{-\beta_{2n+2}-a^\ell_1}{2}<c<\frac{-\beta_{1}-a^\ell_1}{2}\} $.
\end{lemma}


\medskip
Now, let us formulate an inverse problem recovering the relaxation tensor of the q-EBM from the C-ev's of the FOE as follows.

\medskip
\noindent
{\bf Inverse problem:} Recover $\alpha_j,\, \beta_j$ with $1\le j\le 2n+2$ of \eqref{s4.4} and $\lambda_0+2\mu_0$ by knowing sets of the clustered eigenvalues of \eqref{s4.14}.

\medskip
Then, we have the following answer to the above inverse problem.
\begin{Th}\label{s_thm5.2}
By knowing two clusters of eigenvalues associated with $\ell=\ell_1,\,\ell_2\in{\mathbb N}$, we can recover $\alpha_j,\, \beta_j$ with $1\le j\le 2n+2$ of \eqref{s4.4} and $\lambda_0+2\mu_0$. 
\end{Th}
\begin{proof}
In Theorem 1 of \cite{CDDLN2},  we recover $D,\,b_j,\,r_j$, $1\leq j\leq N$ of \eqref{s5.1} by using two clusters of eigenvalues. In the same way, we can recover $\lambda_0+2\mu_0$, $\alpha_j$, $\beta_j$, $1\leq j\leq 2n+2$.
\end{proof}


\begin{remark}
We can find suitable $\mu_0$ which depends on the order of \eqref{s4.2} and then get the associated $\lambda_0$, and $(v_0^j)^2$, $(q_0^j)^2$ for $0\leq j \leq n$ in \eqref{s4.1}. In fact, from the second equation and the fifth equation in \eqref{s4.3}, we have 
\begin{equation}\label{s5.4}
\begin{aligned}
\frac{\alpha_{j+1}}{\beta_{j+1}}=\frac{4\mu_0}{3}(v_0^j)^2,\quad \frac{\alpha_{j+n+2}}{\beta_{j+n+2}}=(\lambda_0+\frac{2}{3}\mu_0)(q_0^j)^2,\qquad 0\leq j \leq n.
\end{aligned}
\end{equation}
Since $ \sum_{j=0}^{n}(v_0^j)^2=1 $ and $ \sum_{j=0}^{n}(q_0^j)^2=1 $, we obtain from \eqref{s5.4} that
\begin{equation}\label{s5.5}
\begin{aligned}
\mu_0=\frac{3}{4}\sum_{j=1}^{n+1}\frac{\alpha_{j}}{\beta_{j}},\quad \lambda_0+2\mu_0=\sum_{j=1}^{2n+2}\frac{\alpha_{j}}{\beta_{j}}.
\end{aligned}
\end{equation}
Note that, by the first equation in \eqref{s5.5}, $\mu_0$ depends on the ordering \eqref{s4.2}.
By using \eqref{s5.5}, we can determine $\lambda_0$ by
\begin{equation}\label{s5.6}
\begin{aligned}
\lambda_0=\sum_{j=1}^{2n+2}\frac{\alpha_{j}}{\beta_{j}}-\frac{3}{2}\sum_{j=1}^{n+1}\frac{\alpha_{j}}{\beta_{j}}=\sum_{j=n+2}^{2n+2}\frac{\alpha_{j}}{\beta_{j}}-\frac{1}{2}\sum_{j=1}^{n+1}\frac{\alpha_{j}}{\beta_{j}}.
\end{aligned}
\end{equation}
From \eqref{s5.4}, we can derive $(v_0^j)^2$ and $(q_0^j)^2$ for $0\leq j \leq n$, respectively. 
\end{remark}
    





\section{Conclusions and discussions }
Being stimulated by the statement given in \cite{YP} on the importance of the EBM for analyzing an inelastic effect coming from the upper mantle to normal modes of the Earth, we analyzed the cluster eigenvalues of the Earth using a three-dimensional homogeneous isotropic EBM. When we started this research, we first noticed that an explicit form of the relaxation tensor was missing for the three-dimensional isotropic generalized Burgers model (EBM') not necessarily homogeneous. Even though in our previous research, we obtained the relaxation tensor for a more general EBM, including the anisotropic case, it lacked explicitness. It was not as easy as we thought to fill this gap. To overcome this difficulty, we found that the volumetric and deviatoric decompositions worked very well. Thus, we developed a new method to derive an explicit form of the relaxation tensor of the aforementioned EBM', which is a fundamental contribution to the people studying EBM'.

On the other hand, for generating the C-ev's of the EBM, we had to pay the price of using the aforementioned decompositions. That is, for the instantaneous part of the q-EBM, we had to find a $u$ that is an eigenfunction for both of the operators $Q_A$ and $Q_B$ with a traction-free boundary condition, simultaneously (see Section 3). We refer to this eigenvalue problem as the simultaneous eigenvalue problem. Comparing the set of eigenvalues for the instantaneous part of the q-EBM with the traction-free boundary condition and that of the simultaneous eigenvalue problem, the latter one is much smaller than the former one. This is a new feature of the C-ev's for the FOE, showing its disadvantage for seeking eigenvalues of the q-EBM. Nevertheless, the C-ev's have a very nice structure which enables to recover the relaxation tensor of the homogeneous EBM. Based on these, we speculate that the C-ev’s have the potential to become popular
in the fields of geophysics and material science.

As for some future works, we provide numerical verifications of the structure of the mentioned C-ev's and the mentioned recovery in our next paper. Also, we are planning to provide an explicit form of the relaxation tensor for the mentioned more general EBM by using the joint spectral decomposition of elasticity tensors in the EBM by assuming that they are commutative. Furthermore, we are planning to use eigenfunctions of the C-ev’s as correlation functions to numerically find the C-ev’s. 



\begin{thebibliography}{99}

\bibitem{CDDLN1}
S. Chen, M. de Hoop, Y. Deng, C-L. Lin, G. Nakamura, 
Clustered eigenvalue problem for glassy state relaxation and its inverse problem, arXiv: 2509.16714.

\bibitem{CDDLN2} 
S. Chen, M. de Hoop, Y. Deng, C-L. Lin, G. Nakamura, Inverse spectral problem for glassy state relaxation approximated by Prony series, arXiv:2512.00507.


\bibitem{CK} {D.~Colton and R.~Kress}, {\it Inverse
Acoustic and Electromagnetic Scattering Theory}, 4th Edition,
Springer-Verlag, Berlin, 2019.

\bibitem{Dafermos}
C. Dafermos, An abstract Volterra equation with application to linear viscoelasticity, J. Differential Equations \, 7 (1970) pp. 554-569.

\bibitem{HKLNT}
M. de Hoop, M. Kimura, C-L. Lin, G. Nakamura and K. Tanuma, Anisotropic extended Burgers model, its relaxation tensor and properties of the associated Boltzmann viscoelastic system,  Siam Journal on Applied Math., {\bf 85} (2025), 1172-1188.

\bibitem{HKLNT_IPMS}
M. de Hoop, M. Kimura, C-L. Lin, G. Nakamura and K. Tanuma, New derivation of relaxation tensor for anisotropice extended Burgers model, Inverse Problems: Modelling and Simulation, Extended Abstracts of the IPMS Conference 2024, edited by A. Hasanov, R. Novikov, K. Bockstal, Birkhauser Cham, 2025, https://doi.org10.1007/978-3-031-87213-6.


%\bibitem{kong}Y. Deng, L. Kong, H. Liu and L. Zhu, {\it Elastostatics with multi-layer metamaterial structures and an algebraic framework for polariton resonances}, ESAIM: M2AN, {\bf 58} (2024) 1413--1440.

%\bibitem{spectraltan}Y. Deng, H. Li and H. Liu, {\it On spectral properties of Neuman-Poincar\'{e} operator and plasmonic resonances in 3D elastostatics}, J. Spectr. Theory, {\bf 9} (2019), 767-789.

%\bibitem{DLL}Y. Deng, H. Li and H. Liu, {\it Spectral properties of Neumann-Poincar\'{e} operator and anomalous localized resonance in elasticity beyond quasi-static limit}, J. Elast, {\bf 140} (2) (2020), 213-242.

\bibitem{Flugge}
W. Flugge, Viscoelasticity, Springer-Verlag, 1975.

\bibitem{Ivins}
E. Ivins, L. Caron, S. Adhikari and E. Larour, Notes on an compressible extended Burgers model of rheology, Geophys. J. Int., 228 (2022) pp. 1975-1991.

\bibitem{KLLNY}
M. Kimura, C-L. Lin, Y. Liu, G. Nakamura, Z. Yang, The derivation of a double-term time-fractional wave equation for the fractional extended Maxwell model, 2026, preprint.

\bibitem{Kor02}
B. G. Korenev, {\it Bessel functions and their applications}, Taylor \& Francis, 2002.

\bibitem{LauFaul_2019}
H. Lau and U. Faul, Anelasticity from Seismic to Tidal Timescales: {{Theory}} and Observations, Physics of the Earth and Planetary Interiors, 508 (2019) pp. 18-29

\bibitem{LS}
P. Loreti, D. Sforza, Viscoelastic aspects of glass relaxation models, Physica A, 526 (2019) 120768.

\bibitem{Harry}
H. Matchette-Downes, Some studies on the computation and interpretation of seismic interface waves and modes in Earth's
mantle (2021) Dissertation submitted to MIT.

 

%https://doi.org/10.1016/j.physa.2019.04.004

\bibitem{Nikabadze}
M. Nikagadze, Eigenvalue problems for tensor-block matrices and their applications to mechanics, J. Math. Sci., 250 (2020) pp. 895-931.

\bibitem{YP}
D. Yuen, W. R. Peltier, Normal modes of the viscoelastic earth.  Geophysical Journal International 69.2 (1982) pp. 495–526. 

\bibitem{Zemanian}
A. Zemanian, Distribution Theory and Transform Analysis, Dover Publications, New York, 1987.
\end{thebibliography}

\begin{figure*}[htb]
  \centering
\fbox{
\includegraphics[width=0.6\textwidth]{CLin2.png}}
\caption*{The extended Burgers model consists of one Maxwell model $M_0$ and $n$ numeber of the Kelvin-Voigt models $K_i$, 
$i=1,\cdots,n.$ Also, they are connected in series. Further, the superscripts $s$ and $d$ refer to springs and dashpots, respectively.}
    \label{fig:EBM}
\end{figure*}




\end{document}
----------------------------------------------
While capturing the full range of the observed mechanical behavior of the upper mantle materials and reconciling the behavior of the mantle across the ``complete'' spectrum of geodynamic time-scales, the EBM has emerged as the model of choice \cite{YuenPeltier_1982}. Dissipation data with frequency dependence of attenuation across, say, six orders of magnitude, are indeed typically best fitted to this model. The EBM yields a consistent description across normal modes, seismic body waves and body tides, and from laboratory ultrasound to the Wilson cycle; see Lau and Faul \cite{LauFaul_2019}. It also captures the long time scale of mantle convection on which Earth behaves viscously. 
% See Irving \textit{et al}. \cite{Irving_2018} for the pursuit of a likewise self consistent model in Earth's outer core. 
The Andrade creep model, and EBM, have been introduced in the astronomical community by Efroimsky and Lainey \cite{EfroimskyLainey_2007}; it has been applied, for example, to Enceladus' dissipation of energy due to forced libration (see Gevorgyan \textit{at al}. \cite{GevorgyanBoueRagazzoRuizCorreia_2020}).





-----------------------------------------------
\bibitem{AndersonMinster_1979}
D. Anderson and J. Minster, The Frequency Dependence of {{Q}} in the {{Earth}} and Implications for Mantle Rheology and {{Chandler}} Wobble, Geophysical Journal International, 58 (1979) pp. 431-440.

\bibitem{Browaeys}
J. Browaeys and S. Chevrot, Decomposition of the elastic tensor and geophysical applications, Geophysical Journal International, 159 (2004) pp. 667-678.

\bibitem{EfroimskyLainey_2007}
M. Efroimsky and V. Lainey, Physics of Bodily Tides in Terrestrial Planets and the Appropriate Scales of Dynamical Evolution, Journal of Geophysical Research: Planets, 112 (2007) pp. 2007JE002908.

\bibitem{GevorgyanBoueRagazzoRuizCorreia_2020}
Y. Gevorgyan, G. Bou{\'e}, C. Ragazzo, L. Ruiz, S. Lucas and A. Correia, Andrade Rheology in Time-Domain. {{Application}} to {{Enceladus}}' Dissipation of Energy Due to Forced Libration, Icarus,  343 (2020) pp. 113610.

\bibitem{Hetland}
E. Hetland and B. Hager, The effects of rheological layering on post-seismic deformation, Geophys. J. Int., 166 (2006) pp. 277-292.

\bibitem{Ivins}
E. Ivins, L. Caron, S. Adhikari and E. Larour, Notes on an compressible extended Burgers model of rheology, Geophys. J. Int., 228 (2022) pp. 1975-1991.

\bibitem{LauFaul_2019}
H. Lau and U. Faul, Anelasticity from Seismic to Tidal Timescales: {{Theory}} and Observations, Physics of the Earth and Planetary Interiors, 508 (2019) pp. 18-29.

\bibitem{LNQ}
H. Li, V. Nistor and Y. Qiao, Uniform Shift Estimates for Transmission Problems and Optimal Rates of Convergence for the Parametric Finite Element Method. In: I. Dimov, I. Farago and L. Vulkov (eds), Numerical Analysis and Its Applications. NAA 2012. Lecture Notes in Computer Science, vol. 8236. Springer, Berlin, Heidelberg. %https://doi.org/10.1007/978-3-642-41515-9_2




-------------------------------------------
\hspace{2.0cm}
  \subfigure[]{
  \includegraphics[width=0.19\textwidth,height=0.17\textwidth]{fig/SLS.png}
  %\label{SLS}
  }

-------------------------------------------\begin{comment}

There is a creep phenomenon that distinguishes viscoelasticity from elasticity. Boltzmann proposed the superposition principle (abbreviated by BSP) to understand this phenomenon. It says that in a linear viscoelastic material subjected to a certain stress history, the strain at a certain time is equal to the sum of the increments of strain caused at that time by each of the previous stress increments acting independently for the entire stress history. Mathematical description of this is given as
\begin{equation}\label{B1.1}
\sigma(x,t)=\int_0^t G(x,t-s)\partial_s e(x,s)\,ds.
\end{equation}
Here $x$ is in a reference domain $\Omega\subset{\mathbb R}^d$ with $d=1,2,3$ for a viscoelastic medium, $t$ is time in the interval $[0,T)$ with $0<T\le \infty$.
Also, $\sigma(x,t)$, $e(x,t)$, $G(x,t)$ are the stress, strain and relaxation tensor, respectively. This is nothing but the stress-strain relation for BVS. Note that the superposition principle does not provide a way to give appropriate assumption on $G(x,t)$. Therefore, modeling people have been using spring-dashpot models (abbreviated by SDs).

\end{comment}
------------------------
\captionsetup{labelsep=none}
\begin{figure}[htbp] 
  \centering
  \subfigure[]{
  \includegraphics[width=0.10\textwidth]{fig/S1.png}
  }
  \hspace{3.0cm}
  \subfigure[]{
  \includegraphics[width=0.10\textwidth]{fig/D1.png}
  }
  \hspace{3.0cm}
  \subfigure[]{
  \includegraphics[width=0.10\textwidth]{fig/M1.png}
  }
  \hspace{3.0cm}
  \subfigure[]{
  \includegraphics[width=0.2\textwidth,height=0.2\textwidth]{fig/KV.png}
  %\label{Kelvin-Voigtモデル}
  }
  \hspace{2.0cm}
 \subfigure[]{
  \includegraphics[width=0.12\textwidth,height=0.2\textwidth]{fig/B.png}
  %\label{B}
  }
  
  \caption{$\,\,\,$S, D, M, KV, B are schematic pictures of a spring, dashpot, Maxwell model, Kelvin-Voigt model, and a Burgers model, respectively}
\end{figure}
---------------


















